% 本文件由 LaTeX 排版 Agent 根据有效 translation.json 自动生成。
% author=Jerome; work_id=3011; title=驳斥白拉奇派

\chapter{驳斥白拉奇派}

% structure: p0001 source=h1 target=omitted reason=page-title-already-rendered-by-parent

卷之一\newline{}卷之二\newline{}卷之三\par

\SourceRecord{Translated by W.H. Fremantle, G. Lewis and W.G. Martley.；Nicene and Post-Nicene Fathers, Second Series；Vol. 6.；Edited by Philip Schaff and Henry Wace.；Buffalo, NY: Christian Literature Publishing Co.,；1893.；<http://www.newadvent.org/fathers/3011.htm>.}

% structure: p0001 source=h1 target=section reason=page-title

\section{驳斥白拉奇派（卷一）}

反白拉奇的对话录是热罗尼莫最后一部论战作品，写于417年，即他去世前三年。该作虽因年迈而略显冗长，却不乏其昔日活力。他对这一主题的审视较先前专论更为冷静，主要因为这多少偏离了他本人的思想轨道。促使他关注此事的，是他对奥斯定日益增长的敬意，以及414年年轻的西班牙人奥罗修从奥斯定处前来投师门下。白拉奇也来到了巴勒斯坦，他的教义在415年由若望主教主持的耶路撒冷小型公会议和416年在狄奥斯波里举行的第二次公会议中受审查后，获准共融。热罗尼莫似乎未参与这些进程，且与若望主教和睦相处近二十年，无疑不愿反对他。但他已视白拉奇为沾染了异端式的\Quote{不敬}（第一卷第28章），认为这远甚于道德之恶；并从他致特斯丰的信（第133号）中可见，他将白拉奇与奥力振主义和鲁菲诺联系起来；他运用其广博的圣经知识来应对这场争论。他引用了白拉奇的一部著作，尽管只给出了标题和章数，直至第100章（第一卷第26至32章）；而且，尽管他的信念有时显得薄弱，且某些段落（第一卷第5章，第二卷第6至30章，第三卷第1章）让人有机会注意到他实际上（即使是无意识地）倾向于白拉奇的观点，并称他为\Quote{协同论者}，而非像奥斯定那样是彻底的预定论者，但该对话录整体上思路清晰，为我们的知识做出了实质贡献。虽然其语气不如他的苦修专论激烈，但它似乎激起了对他最强烈的敌意。白拉奇的支持者袭击并焚烧了白冷的隐修院，热罗尼莫本人只得以躲进一座塔楼才得以逃脱。他的苦难以及教宗依诺增爵对他的干预，可参见附有导言注释的第135至137号书信。\par

以下是论证概要：奥斯定派的亚提古立即（第1章）提出问题：你是否肯定，如白拉奇所肯定的那样，人可以不犯罪而生活？白拉奇派的克里托布鲁答道：是的，但我不加上如人们所指控我们的\Quote{没有天主的恩宠}。事实上，我们拥有自由意志正是出于恩宠。亚提古回答说，但那恩宠是什么？是仅仅我们原初的本性，还是每项行动都需要？克里托布鲁答道，每项行动都需要（第2章）；然而，人们几乎不会说，没有恩宠我们就不能修好一支笔（第3章），因为若是如此，我们的自由意志又在何方？但亚提古说（第5章），圣经论及我们在一切事上都需要天主的援助。在那种情况下，克里托布鲁说，所应许的赏报就必须归给天主，而非归给我们，因为是天主在我们内工作。随后回到所陈述的第一点，亚提古问道，能不犯罪的可能性是指个别行动，还是指向整个一生？回答当然是整体与部分都包括。但我们希望或愿意不犯罪；为何实际上我们却不犯罪？因为（第8章）我们没有完全运用意志。但（第9章）从未有人能一生不犯罪。然而，白拉奇派说，天主命令我们成为成全的，而祂不会命令不可能的事。约伯、匝加利亚和依撒伯尔被描述为完全正义的。答说：不然（第12章），他们每人都被指称有缺点。若望说：\BibleQuote{由天主所生的人不犯罪}（\BibleRef{若一 3:9}）；然而，\Quote{若我们说我们没有罪，便是自欺}（\BibleRef{若一 1:8}）。宗徒们虽被命令成为成全的（第14章），却并未成全；圣保禄说：\BibleQuote{我不以为自己已经取得了}（\BibleRef{斐 3:13}）。称人为正义和成全，只是与他人相比（第16章），或是因为普遍服从天主的旨意（第18章），或是根据其特殊特性（第19章），正如我们可以说一位主教在其职位上卓越，尽管他可能并未达到\WorkTitle{牧函}的理想（第22章）。\par

讨论现在转向白拉奇著作中的言辞。\Quote{所有人都受自己的意志支配}（第27章）。不；因为基督说\BibleQuote{我来不是为行我的旨意}（\BibleRef{若 6:38}）。\Quote{在审判中，恶人必不得存留}。但我们必须区分将被消灭的不敬之人或异端者（第28章）与将被宽恕的基督徒罪人。他的一些说法相互矛盾或流于琐碎（第29、30章）。\Quote{旧约中已应许了天国}。是的，但新约中更为全面。回到第一个论点：\Quote{人若愿意便能不犯罪}，白拉奇派说：如果像自发产生而无结果的情欲这类事被认为是有罪的，我们就将罪归咎于造物主；对此仅回答：尽管我们不能理解天主的道路，却不可质疑祂的正义。在本书其余部分，亚提古单独发言，通览旧约，指明每位圣人都会犯下某种罪，罪虽出于无知或半意识，仍带来定罪。\par

% structure: p0005 source=h2 target=subsection reason=relative-heading-rank; base=h2

\subsection{序言}

1. 在致克特西丰的书信（其中我已答复所提出的问题）之后，弟兄们屡次向我抱怨，想知道为何我迟迟不履行承诺的著作，即我应允答复所有宣扬无情之人的狡诈。因为人人都知道斯多葛派和逍遥学派（即旧学园）的争论：其中一些人断言，\OriginalTerm{情感}{\GreekText{πάθη}}（可称为情绪，如悲伤、喜悦、希望、恐惧）可以从人心中彻底根除；另一些人则认为，它们的力量可以被削弱，可以被控制和约束，就像桀骜不驯的马匹用特殊嚼子勒住一样。西塞罗在\WorkTitle{图斯库鲁姆论辩录}中解释了他们的观点，而奥力振在其\WorkTitle{杂论}中试图将它们与教会真理相融合。我暂且不提摩尼、普里西利安、伊博拉的厄瓦格利乌斯、约维尼安，以及几乎遍布叙利亚全境的异端者——他们因名字的曲解而通常被称为\OriginalTerm{玛撒里安派}{Massalians}，希腊语为\OriginalTerm{尤希特派}{Euchites}——所有人都认为，人的德行和知识可以达到完美，并达到（我不说相似于，而是）等同于天主的地步；他们甚至断言，一旦达到完美的顶峰，连思想和无知之罪对他们也是不可能的。虽然在我先前写给克特西丰并针对他们错误的信中，我趁时间允许已触及几点，在我现在努力雕琢的这本书中，我将遵循苏格拉底的方法。将陈述双方能说的一切；当双方表达观点时，真理将会显明。奥力振持独特立场：一方面认为人性不可能一生无罪，另一方面又认为人转向更美善之事时，可以变得如此坚强以至于不再犯罪。\newline{}\par

2. 我还要对那些说我因嫉妒而写作此书的人说几句话。我从未宽容过异端者，并尽力使教会的仇敌成为我自己的仇敌。赫尔维迪乌斯曾反对圣玛利亚的终身童贞。我回答他——我从未在肉身上见过他——难道是出于嫉妒吗？约维尼安，他的异端现正煽风点火，并在我不在时扰乱了罗马的信德，他口才如此贫乏，文风如此有害，与其说是嫉妒的对象，不如说是怜悯的对象。我尽我所能也回答了他。鲁菲努斯竭尽全力传播奥力振的亵渎言论和\WorkTitle{论首要原理}（\OriginalTerm{论首要原理}{\GreekText{Περὶ} \GreekText{᾿Αρχῶν}}），不是在一个城市，而是在全世界。他甚至以殉道者庞非吕的名义出版了欧瑟伯的\WorkTitle{为奥力振辩护}的第一卷，并且好像奥力振说得还不够似的，又为他说出了一卷新的著作。我回答他，就要被指责为嫉妒吗？他的口才难道是一条汹涌的急流，使我因恐惧而不敢写或口授任何东西作为答复吗？帕拉迪乌斯，像恶奴一样，试图为同一个异端注入活力，并因我翻译希伯来文而激起对我的新偏见。我难道嫉妒如此杰出的才智和贵胄吗？即使现在，\BibleRef{得后 2:7}不法的奥秘已在运行，人人都在谈论自己的观点：然而，似乎只有我一人因嫉妒所有其他人的荣耀而充满嫉妒；我是如此可怜，竟嫉妒那些不值得嫉妒的人。所以，为了向所有人证明我不是恨人而是恨他们的错误，并且不希望诽谤任何人，而是哀叹那些被虚假知识所欺骗之人的不幸，我用了阿提库斯和克里托布鲁斯的名字，来表达我们自己和反对者的观点。事实上，所有我们持守大公信仰的人，都希望并渴望，当异端被谴责时，这些人能够得到改造。无论如何，如果他们执迷不悟，过错不在我们这些写作的人，而在他们，因为他们宁愿选择谎言而不选择真理。对我们的诽谤者（他们的咒骂落在自己头上）有一个简短的回答：摩尼教的教义谴责人的本性，摧毁自由意志，并且取消了天主的帮助。再说，人自称是唯独天主所是的东西，那是明显的疯狂。让我们沿着王道行走，既不偏右也不偏左；并始终相信，我们意愿的热切是由天主的帮助所管束的。若有人嚷着说他被诽谤，并夸口说他和我们想法一致；那么只有当他公开而真诚地谴责相反的观点时，他才显出他赞同真信德。否则，他的情况就如先知所描述的：\BibleRef{耶 3:10}\Quote{虽然如此，她那奸诈的姊妹犹大还没有全心全意归向我，只是假意。}若跟随你以为是善的恶，这是较小的罪；不敢捍卫你确知是善的恶，这是更大的罪。如果我们不能忍受威胁、不义、贫穷，又如何能战胜巴比伦的火焰？不要因虚假的和平而失去我们靠战争保存的东西。我不愿让恐惧教我背信，而基督已把真信德放在我选择的能力之内。\newline{}\par

% structure: p0008 source=h2 target=omitted reason=duplicate-page-title

\Person{阿提库斯}：\Quote{我听说，克里托布鲁斯，你写道：人若愿意，可以无罪；并且天主的诫命是容易的。请告诉我，这是真的吗？}\newline{}\par

\Person{克里托布鲁斯}：\Quote{阿提库斯，是真的；但我们的对手并未按我所指的意义来理解这些话。}\newline{}\par

\Person{阿提库斯}：\Quote{那么它们是否如此含混，以致在意义上产生了分歧？我不要求同时回答两个问题。你提出了两个命题：一是人若愿意可以无罪；二是天主的诫命容易。因此，虽然它们是一起说出的，但让它们分别讨论，这样，当我们的信仰看来是一致时，就不会因相互误解而产生争吵。}\newline{}\par

\Person{克里托布路斯}说：阿提库斯，我说过人能够无罪，如果他愿意；并非像某些人恶意曲解我们所说的，没有天主的恩宠（这种想法本身就是不敬），而仅仅是他能够，如果他愿意；但预先假定天主恩宠的帮助。\par

\Person{阿提库斯}说：那么，天主是你恶行的作者吗？\par

\Person{克里托布路斯}说：绝非如此。但若我内有任何善，都是通过祂的推动和协助而达至完善的。\par

\Person{阿提库斯}说：我的问题不是关于自然构造，而是关于行动。因为谁怀疑天主是万物的创造者呢？我希望你告诉我这一点：你所行的善，是你的，还是天主的？\par

\Person{克里托布路斯}说：是我的，也是天主的：我工作，祂协助。\par

\Person{阿提库斯}说：那么，为何人人都认为你废除了天主的恩宠，并主张我们的一切行动都出于我们自己的意志？\par

\Person{克里托布路斯}说：阿提库斯，我很惊讶你问我别人错误的缘由，想知道我没有写的东西，而我写的东西非常清楚。我说过人能够无罪，如果他愿意。我加上了\Quote{\Emph{没有天主的恩宠？}}吗？\par

\Person{阿提库斯}说：没有；但你没有添加任何话这一事实暗示了你否认恩宠的必要。\par

\Person{克里托布路斯}说：不，相反，我没有否认恩宠这一事实，应被视为等同于肯定它。认为我们未断言的就是否认，这是不公平的。\par

\Person{阿提库斯}说：那么，你承认，人能够无罪，如果他愿意，但要有天主的恩宠。\par

\Person{克里托布路斯}说：我不仅承认，而且自由地宣扬它。\par

\Person{阿提库斯}说：那么，废除天主恩宠的人就是错误的。\par

\Person{克里托布路斯}说：正是这样。或者更确切地说，他应被认为是不敬的，因为万物都由天主的旨意管理，我们的存在、个别选择和欲望的能力都归功于造物主天主的良善。因为我们有自由意志，并根据我们自己的选择倾向善或恶，这是祂恩宠的一部分，祂按自己的肖像和模样造了我们。\par

2. \Person{阿提库斯}说：没有人怀疑，克里托布路斯，万物都取决于万有创造者的判断，我们所有的一切都应归于祂的良善。但我希望知道，针对这种你归功于天主恩宠的能力，你是将其视为我们受造时赐予的恩赐的一部分，还是认为它在我们的个别行动中是活跃的，以致我们持续利用它的协助？或者情况是，我们被一次创造并赋予了自由意志后，凭自己的选择或力量做我们选择的事？因为我知道，你们一派中有很多人将万物归于天主恩宠，其意义是，他们理解意志的力量不是特殊的恩赐，而是一般的恩赐，也就是说，不是在每个个别时刻赐予的，而是在创造时一次性赐予的。\par

\Person{克里托布路斯}说：并非如你所说；但我坚持两种立场：我们是因天主的恩宠被造成现在的样子，并且在我们的各项行动中，我们获得祂的协助的支持。\par

\Person{阿提库斯}说：那么，我们达成一致：在善工中，除了我们自己的选择能力，我们还依靠天主的帮助；在恶工中，我们受魔鬼的唆使。\par

\Person{克里托布路斯}说：正是如此；在这一点上没有意见分歧。\par

\Person{阿提库斯}说：那么，那些剥夺我们在个别行动中天主的帮助的人是错误的。圣咏作者歌唱说：\BibleQuote{除非主建造房屋，建造的人徒然劳苦；除非主看守城池，守醒的人徒然警醒；}还有其他类似的章节。但这些人试图通过乖僻的，或者更确切说是可笑的解释，将他的话曲解成别的意思。\par

3. \Person{克里托布路斯}说：当你有我自己的回答时，我有义务反驳别人吗？\par

\Person{阿提库斯}说：你的回答是何种意思？说他们是对的，还是错的？\par

\Person{克里托布路斯}说：有什么必要迫使我与别人的意见对立？\par

\Person{阿提库斯}说：你受讨论规则和尊重真理的约束。难道你不知道每个断言要么肯定，要么否定，而被肯定或否定的事物应被视为或好或坏？因此，你必须承认——虽然并非出于你的本意——我的问题所涉及的说法要么是好事，要么是坏事。\par

\Person{克里托布路斯}说：如果在个别行动中我们必须有天主的帮助，那么是否意味着我们不能制作一支笔，或修好它？我们不能在没有天主协助的情况下写字、沉默或说话、坐、站、走或跑、吃或禁食、哭或笑等等吗？\par

\Person{阿提库斯}说：从我的观点看，这显然是不可能的。\par

\Person{克里托布路斯}说：那么，如果我们连这些事都不能没有天主而做，我们如何拥有自由意志，如何能保守天主对我们的恩宠？\par

4. \Person{阿提库斯}说：自由意志恩宠的赐予，并非如此以致废除了天主在特定行动中的支持。\par

\Person{克里托布路斯}说：天主的帮助并非变得毫无意义；因为受造物是通过一次赐予他们的自由意志的恩宠而得以保存的。因为如果没有天主，除非祂在每一行动中协助我，我什么都不能做。祂既不能公义地为我的善行加冕，也不能惩罚我的恶行，而是在每种情况下，祂要么收回祂自己的，要么谴责祂所给予的协助者。\par

\Person{阿提库斯}说：那么，坦率地告诉我，你为什么废除天主的恩宠。因为你在部分中摧毁的任何东西，你必然在整体中否认它。\par

\Person{克里托布路斯}说：当我断言我是被天主如此创造的，以致因天主的恩宠，做一件事或不做的能力被置于我选择的能力之内时，我并没有否认恩宠。\par

\Person{阿提库斯}说：那么，一旦自由意志的能力被赐予，天主就在我们的善行上睡着了；我们不需要祈求祂在个别行动中协助我们，因为做一件事（如果我们选择）或不做的能力取决于我们自己的选择和意志。\par

5. \Person{克里托布路斯}说：如同其他受造物的情况一样，创造的既定条件被遵守；因此，一旦自由意志的能力被赐予，一切都被交给了我们自己的选择。\par

\Person{阿提库斯}说：由此得出，如我所说，我不应在行为的细节上祈求天主的协助，因为我认为它是一次赐予的。\par

\Person{克里托布路斯}说：如果祂在一切事上与我合作，结果就不再是我的，而是那位协助者的，或者更确切地说，是在我内并同我工作的那位的；更是因为我离了祂什么也不能做。\par

A. \Quote{你们没有念过\BibleRef{罗 9:16}上记载的：\Quote{这不在乎那愿意的，也不在乎那奔跑的，只在乎那显示仁慈的天主}吗？}由此我们明白：愿意和奔跑是属于我们的，但将我们的愿意和奔跑付诸实行，则属于天主的仁慈，并且如此成就，以致一方面在愿意和奔跑中保存了\OriginalTerm{自由意志}{free will}，另一方面，在完成我们的愿意和奔跑时，一切都归于天主的权能。当然，我现在应当援引圣经多次的见证，以表明圣人们在行为的细节中恳求天主的助佑，并在他们的各种行动中渴望祂作他们的帮助者和保护者。请通读《圣咏集》和圣人们的一切言论，你会发现他们的行动从未缺少对天主的祈祷。这清楚地证明：要么你否认恩宠，并将它从生命中驱逐出去；要么，如果你承认它在各个部分中存在——这显然很违背你的意愿——你就必须转向我们这些为人类保存自由意志，但如此限制自由意志以致不否认天主在每一行动中的助佑之人的观点。\par

6. C. 这是一个诡辩的结论，仅仅是逻辑技巧的卖弄。没有人能剥夺我的自由意志；否则，如果天主真是我行事中的帮助者，那报酬就不应归于我，而应归于那在我内工作的那一位。\par

A. 你尽管利用你的自由意志吧；用你的舌头武装起来反对天主，如果愿意，你就在其中证明自己是自由的，可以亵渎。但更进一步，你的观点已毋庸置疑，你职业的幻象已昭然若揭。现在，让我们回到我们讨论的起点。你刚才说，假定有天主的助佑，如果人愿意，就可以无罪。请告诉我，这能持续多久？永远，还是仅仅短暂的时间？\par

C. 你的问题是不必要的。如果我说短暂时间，永远也同样被包含在内。因为无论你允许短暂时间，你都会承认它可以永远持续。\par

A. 我不太明白你的意思。\par

C. 你如此迟钝，连明显的事实都认不出来吗？\par

7. A. 我不以我的无知为耻。双方应当对争论主题的定义达成一致。\par

C. 我主张：一个人如果一天能避免犯罪，那么另一天也能；如果能两天，就能三天；如果能三天，就能三十天；如此类推，三百天、三千天，或随他愿意多久。\par

A. 那么就立刻说，如果人愿意，他可以永远无罪。我们能做任何我们喜欢的事吗？\par

C. 当然不能，因为我不能做所有我喜欢的事；但我只说，如果人愿意，他可以无罪。\par

A. 请告诉我：你认为我是人还是兽？\par

C. 如果我怀疑你是人还是兽，我就承认自己是兽了。\par

A. 那么，如你所说，既然我是人，为什么当我愿意并热切渴望不犯罪时，却仍犯罪呢？\par

C. 因为你的选择不完全。如果你真的不愿意犯罪，你就真的不会犯罪。\par

A. 那么，指责我没有真正渴望的你，是否因为你有真正渴望而无罪呢？\par

C. 好像我是在谈论我自己——我承认是罪人——而不是谈论那些决意不犯罪的少数例外之人（如果有的话）。\par

8. A. 尽管如此，我这个提问者和你这回答者都认为自己是罪人。\par

C. 但如果我们愿意，我们能够不是这样。\par

A. 我说我不愿意犯罪，无疑你也有同感。那么为什么我们两人所愿意的，却都不能做到呢？\par

C. 因为我们不愿意得不够完全。\par

A. 请给我指出我们哪位祖先有完全的意愿并完全地拥有能力？\par

C. 那并不容易。当我说如果人愿意，他可以无罪时，我并非主张曾经有过这样的人；我只维护抽象的可能性——如果他愿意。因为\Emph{实现的可能性}是一回事，希腊文用\OriginalTerm{可能性}{\GreekText{τῇ} \GreekText{δυνάμει}}表示；\Emph{实现}是另一回事，其对应词是\OriginalTerm{实现}{\GreekText{τῇ} \GreekText{ἐνεργείᾳ}}（参看\BibleRef{创 17:17}、\BibleRef{创 18:12}）。\Person{亚巴郎}和\Person{撒辣}嗤笑那恩许。\par

\BibleRef{创 37:35}。\Person{雅各伯}的过度悲痛。\par

\BibleRef{出 21:12,13}。无意杀人的罪过。\par

\BibleRef{肋 4:2,27}。为无知之罪献祭。\par

\BibleRef{肋 5:3}。为礼仪上的不洁献祭。\par

\BibleRef{肋 9:1}。\Person{亚郎}就职时的献祭。\par

\BibleRef{肋 12:6}。妇女生育后的献祭。\par

\BibleRef{肋 14:1,6}、\BibleRef{肋 16:6}、\BibleRef{肋 12:7}。为癞病人献祭。\par

\BibleRef{肋 15:31}、\BibleRef{肋 16:2,5}。赎罪日为百姓献祭。\par

\BibleRef{肋 22:14}。无知吃圣物；对比\BibleRef{格前 11:27,28}，关乎疏忽领圣事。\par

\BibleRef{户 6:1}。为纳齐尔人献祭。\par

\BibleRef{户 14:7}、\BibleRef{户 7:28,29}。为恳求天主的仁慈献祭。\par

\BibleRef{户 28:15,22}、\BibleRef{户 29:5}、\BibleRef{户 5:11,17}。节期献祭。\par

\BibleRef{户 35:13}。为杀人犯设立的避难城。\par

\BibleRef{申 9:6}、\BibleRef{申 18:13}。警告以色列不可夸耀自己的义。\par

\BibleRef{申 18:9-12}、\BibleRef{申 5:14,15}。完美仅仅用于避免偶像崇拜。\par

\BibleRef{申 22:8}。房顶没有栏杆使人有罪。\par

\BibleRef{申 23:2}。因无意中个人行为而沾染。\par

\BibleRef{苏 7:12}。因\Person{阿干}的罪，百姓成为有罪。\par

\BibleRef{苏 11:19,20}。客纳罕人的种族之罪。\par

\BibleRef{撒上 14:27}。\Person{约纳堂}因尝蜂蜜而成为有罪。\par

\BibleRef{撒上 16:6}。上主看人心，不看外貌。\par

\BibleRef{撒下 4:11}。\Person{依市巴耳}被称为义人。\par

\BibleRef{撒下 6:7,8}。\Person{乌匝}因疏忽被击打。\par

\BibleRef{撒下 24:10}。\Person{达味}统计百姓。\par

\BibleRef{列上 8:46}。\Person{撒罗满}的祈祷：没有不犯罪的人。\par

\BibleRef{列上 14:5}。先知识破\Person{雅洛贝罕}妻子的动机。\par

\BibleRef{列下 4:27}。\Person{厄里亚}看见叔能妇人的心。\par

\BibleRef{编上 2:32}。七十士译本：半先知。\par

\BibleRef{哈 3:1}。武加大译本：\Quote{为无知之罪}的祈祷（\Quote{关于\WorkTitle{示约尼特}}），被认为是承认\BibleRef{哈 1:2-4}中的过分大胆。\par

\BibleRef{则 46:20}。\Person{厄则克耳}重建圣殿的祭献。\par

\BibleRef{耶 10:23}。人的道路不在自己手中。\par

\BibleRef{耶 17:9}。心诡诈。\par

\BibleRef{箴 14:12}。有一条路，人以为正。\par

\BibleRef{箴 19:21}。人心多有计谋。\par

\BibleRef{箴 20:9}。谁能说，我洁净了我的心？\par

\BibleRef{箴 20:17}。谁能夸口自己是洁净的？\par

\BibleRef{训 7:16}。人的心充满邪恶。\par

\SourceRecord{Translated by W.H. Fremantle, G. Lewis and W.G. Martley.；Nicene and Post-Nicene Fathers, Second Series；Vol. 6.；Edited by Philip Schaff and Henry Wace.；Buffalo, NY: Christian Literature Publishing Co.,；1893.；<http://www.newadvent.org/fathers/30111.htm>.}

% structure: p0001 source=h1 target=section reason=page-title

\section{驳斥伯拉纠派（卷二）}

本书很难说是一篇对话的一部分。它更是一个出自圣经的论证，用以证明奥思定派的辩论者\Person{阿提库斯}的论点。从第四章开始，它如同第一卷的最后五章，由一连串取自新约和先知书的圣经经文组成，以显示罪的普遍性，从而驳斥伯拉纠派的断言：若一个人愿意，他可以无罪。因此，像之前一样，我们将给出经文列表及其首词，仅当\Person{热罗尼莫}提出他自己的一些原创评论或值得注意的注解时才引用他的原话。\par

伯拉纠派一开始便重申两难困境：如果诫命是给人遵守的，那么人就可以无罪；如果他因受造而必然成为罪人，那么天主而非他才是罪的作者。至于为无知之罪而规定的祭献，\Person{伯拉纠}则以从旧约到新约的诉求来回应，这引出了对\Person{圣保禄}在\BibleRef{罗 7}中关于与罪斗争的论述的讨论。保禄被认为不是作为罪人说话，而是作为\Emph{人}，从而承认了人性的有罪性。人可能没有根深蒂固的恶习，却不可能没有罪。这导致奥思定派的\Person{阿提库斯}继续列出他的见证清单，说明虽然有人被发现在拒绝恶事上是正义的（\OriginalTerm{邪恶}{\GreekText{κακία}}），却没有人是\OriginalTerm{无罪}{\GreekText{ἀναμάρτητος}}。\par

在\BibleRef{咏 31:5}，一位自称为\Quote{圣洁}的人，却承认自己的过犯。\par

\BibleRef{箴 24:16}解释这一点：\BibleQuote{义人跌倒，却又犯罪。}\par

\BibleRef{箴 18:17}（七十贤士译本与拉丁通行本）：义人开始说话时，便控告自己。\par

\BibleRef{咏 57:3}。罪人一出母胎就与主疏远；也就是说，要么如\Person{圣保禄}在\BibleRef{罗 5:14}所说，他们\BibleQuote{在相似\Person{亚当}的过犯中}犯罪；要么如同\BibleQuote{基督作为首生子，开了童贞母胎}（\BibleRef{出 13:2}）。异端者拒绝承认这个奥迹，这由圣殿东门所预表（\BibleRef{则 44:2}），当大司祭通过之后，那门便又关闭了。\par

\BibleRef{约 4:17}：\BibleQuote{必死的人在天主面前能成义吗？}\par

\BibleRef{约 7:1}：\BibleQuote{人的生命是试炼。}\par

\BibleRef{约 20:21}：\BibleQuote{我若犯了罪，我能做什么？}\par

\BibleRef{约 9:15-16}：\BibleQuote{我若是义人，他也不会听我。}\par

\BibleRef{约 9:29-31}：\BibleQuote{我若用雪水洗身，等等。}\par

\BibleRef{约 10:15}：\BibleQuote{我若为义，等等。}\par

\BibleRef{约 14:4-5}：\BibleQuote{谁能脱离污秽？没有一人。}\par

\BibleRef{箴 16:26}（七十贤士译本）：\BibleQuote{人劳苦在忧愁中。}\par

\BibleRef{约 40:4}：\BibleQuote{我能回答你什么？}\par

\BibleRef{箴 20:9}：\BibleQuote{谁能夸口说自己有一颗清洁的心？}这至少表明诫命并非像\Person{伯拉纠}所说的那么容易。\par

\BibleRef{若一 5:3}：\BibleQuote{他的诫命不是难守的}，以及\par

\BibleRef{玛 11:30}：\BibleQuote{我的轭是容易的}，这些话只有在与犹太教相比较时才为真，并应与\par

\BibleRef{宗 15:10}：\BibleQuote{一个轭……我们的祖先和我们都不能负的}。\par

\BibleRef{雅 4:11}：\BibleQuote{你们论断法律}，也就是说，如果你说对无知之罪的定罪是不合理的。我们都在这些事上犯罪，从\par

\BibleRef{雅 1:20}：\BibleQuote{人的愤怒并不成就天主的正义。}但愤怒经常被谴责，如\par

\BibleRef{箴 15:1}（七十贤士译本）：\BibleQuote{愤怒甚至摧毁智慧人。}\par

\BibleRef{弗 4:26}：\BibleQuote{不要让太阳在你们的愤怒上落下。}\par

\BibleRef{玛 5:22}：\BibleQuote{凡发怒的……难免受议会的审判。}\par

\BibleRef{训 11:19}：\BibleQuote{我是众人中最愚拙的。}这是\Person{基督}以人性身份说的。所以\par

\BibleRef{咏 68:5}：\BibleQuote{天主，你知道我的愚昧。}但是\par

\BibleRef{格前 1:25}：\BibleQuote{天主的愚拙比人更智慧。}\par

\BibleRef{德 1:18}：\BibleQuote{多智慧多忧愁，}显示智慧人对不完美的感受。所以\par

\BibleRef{德 8:7}：\BibleQuote{我憎恨我的生命，}以及\par

\BibleRef{德 8:14}：\BibleQuote{有义人遭遇恶人应得的报应；}也就是说，天主在我们看不见的地方看见邪恶。\par

\BibleRef{德 8:17}：\BibleQuote{人无论怎样劳苦，终不能寻得它；}以及\par

\BibleRef{德 9:2-3}：\BibleQuote{众人遭遇相同。人心……充满邪恶。}\par

\BibleRef{德 10:1}：\BibleQuote{死苍蝇使膏油发臭；}也就是说，几乎人人都被异端或其他过犯玷污。\par

\BibleRef{伯前 2:17-18}：\BibleQuote{审判必须从天主的家开始。}\par

6. 有四种情感搅扰人类，两种关乎现在，两种关乎将来；两种关乎善，两种关乎恶。有忧愁，希腊文称为\OriginalTerm{忧愁}{\GreekText{λύπη}}，与喜乐，希腊文称\OriginalTerm{喜乐}{\GreekText{χαρά}}或\OriginalTerm{快乐}{\GreekText{ἡ} \GreekText{δονή}}，尽管许多人将后一词译为\Emph{voluptas}，快乐；其一归于恶，另一归于善。我们若因不当之事而喜乐，例如财富、权力、荣誉、敌人的厄运或死亡，那就太过分了；反之，若因现世之恶——逆境、流放、贫穷、软弱及亲属之死——而忧伤受折磨，亦为过分，这些都为宗徒所禁止。再者，若我们贪恋我们视为善的事物——继承、荣誉、恒常的繁荣、身体健康等——并在拥有它们时喜乐享受；或我们惧怕我们视为有害的事物。按照斯多葛学派，即\Person{芝诺}与\Person{克里西普斯}，一个成全之人可能摆脱这些情感；按逍遥学派，则困难甚至不可能，而这意见为全部圣经所一贯支持。因此，马加比史家\Person{约瑟夫斯}说，情感可以被制服和管理，而非根除；\Person{西塞罗}的五卷\WorkTitle{图斯库路姆论辩集}亦充满这些讨论。\BibleRef{弗 6:12}按宗徒所说，身体的软弱与天上属灵邪恶的势力攻击我们。同一作者\BibleRef{迦 5:19}告诉我们，肉体的作为与圣神的作为是显而易见的，且彼此相敌，致使我们不做我们所愿意的。若我们不做所愿，反做所不愿，你怎能说人若愿意就能无罪？你看见，无论宗徒或任何信徒，都不能行他所愿的。\BibleRef{伯前 4:8}\BibleQuote{爱遮盖许多罪过}，与其说是过去的罪，不如说是现在的罪，好使我们不再犯罪，只要天主的爱在我们内。因此论及那犯罪的妇人说：\BibleRef{路 7:47}\BibleQuote{她的罪，许多的都赦免了，因为她爱得多。}这告诉我们，做我们所愿的，不仅靠我们自己的力量，也靠天主慈爱对我们意志的援助。\par

7. 现在继续引用圣经：\par

在\BibleRef{若一 1:5}、\BibleRef{若 1:7-8}、\BibleRef{玛 5:14}中，基督和宗徒被称为世界的光。因此世界是黑暗的。\par

\BibleRef{弟前 6:16}。只有天主有永生，且是\Quote{唯一智慧者}；然而其他人，如\Person{提洛王子}\BibleRef{则 28:3}，是分有智慧。同样，我们是纯洁的，但只靠恩宠。因此\par

\BibleRef{若一 1:7}。基督的血洁净我们。\par

\BibleRef{约 25:5-6}。星辰在他眼中也不算纯洁。\par

\BibleRef{迦 2:16}。\BibleQuote{法律不能使任何人成义}，但是\par

\BibleRef{罗 3:1}、\BibleRef{罗 3:24}、\BibleRef{罗 3:28}、\BibleRef{罗 3:30}。因他的恩宠自由地成义，等等。\par

\BibleRef{罗 6:14}。不在法律之下，而在恩宠之下。\par

\BibleRef{罗 9:16}。不在乎那愿意的，而在乎那施仁慈的天主。\par

\BibleRef{罗 9:30-32}。外邦人……因信德获得了义德。\par

\BibleRef{罗 10:2}。基督是法律的终向，使一切相信的人获得。\par

8. 宗徒承认自己工作需要这恩宠。\par

\BibleRef{格前 1:1-3}。恩宠从天主归于你们。\par

\BibleRef{格前 1:7-8}。使你们在恩赐上没有一样不及别人——好使凡有血肉的，在他面前一个也不能自夸。\par

\BibleRef{格前 3:6-10}。保禄栽种了……但天主使生长。\par

\BibleRef{格前 3:18-19}。若有人自以为有智慧，让他变为愚拙。\par

\BibleRef{格前 4:4}。我虽不觉得自己有错，却不能因此成义。\par

\BibleRef{格前 4:7}。你有什么不是领受的呢？\par

\BibleRef{格前 4:19}。我必到你们那里去，\Emph{若主愿意}。\par

9. 宗徒也表明自己需要恩宠。\par

\BibleRef{格前 15:9-10}。靠着天主的恩宠，我成了今日的我，等等。\par

\BibleRef{格后 3:4-6}。我们的充足是从天主来的。\par

\BibleRef{迦 2:16}。我们信了，为要因信德成义。\par

\BibleRef{迦 2:21}。若义德来自法律，基督就白白死了。\par

\BibleRef{迦 3:10}、\BibleRef{迦 3:13}。基督赎我们脱离法律的诅咒。\par

\BibleRef{迦 3:24}。法律是我们的教师，引我们到基督那里。\par

\BibleRef{迦 5:4}。你们这要靠法律成义的，是与基督隔绝了。\par

10.\par

\BibleRef{斐 2:13}。是天主在你们内工作。\par

\BibleRef{得后 3:3}。主是信实的，他必坚固你们。\par

\BibleRef{弟前 6:20-21}。弟茂德啊，要持守所交付你的。\par

\BibleRef{铎 3:4-7}。我们救主天主的慈爱与怜悯救了我们。\par

11. 现在我们转向福音书\Quote{并以基督之灯的亮光补充宗徒之光的闪烁火焰。}\par

\BibleRef{玛 5:22}。\BibleQuote{凡向弟兄动怒的……难免公会的审判。} 我们中间谁在此不被定罪？\par

\BibleRef{玛 5:23-24}。\BibleQuote{先去同弟兄和好。} 有谁觉得这条诫命容易呢？\par

\BibleRef{玛 5:37}。\BibleQuote{你们的话，是就说是，不是就说不是。} 谁曾遵守这条诫命？圣咏作者说\BibleRef{咏 115:11}：所有人都是说谎者。\par

12.\par

\BibleRef{玛 6:34}。\BibleQuote{不要为明天忧虑。} 你做到了吗？\par

\BibleRef{玛 7:14}。\BibleQuote{通向生命的门是窄的。} 你怎能说诫命容易呢？\par

\BibleRef{路 9:58}。\BibleQuote{人子没有枕头的地方。} 这经解释于\par

\BibleRef{依 28:12}。\BibleQuote{接待那疲倦的，这是我的安息。} 并且\par

\BibleRef{依 66:1-2}。\BibleQuote{我安息在谁身上？岂不是在谦卑人身上？} 基督找不着几个可安息的人。那么，怎可说他的诫命容易呢？\par

\BibleRef{玛 9:12-13}。\BibleQuote{我来不是召义人。} \BibleQuote{健康的人不需要医生。} 若世界不是充满罪恶，基督就不会来。因此\par

\BibleRef{咏 12:1}。上主，求你帮助，因为虔诚人绝迹了。\par

\BibleRef{咏 14:1}、\BibleRef{咏 14:3}。他们都败坏了……没有一个行善的。\par

\BibleRef{玛 10:9}。\BibleQuote{你们不要带金……也不要穿鞋。} 谁做到了？连宗徒也没有，因为\par

\BibleRef{宗 12:8}。天使吩咐伯多禄束上鞋。\par

13.\par

\BibleRef{玛 10:22-34}。描述基督跟从者所受的逼迫，并命令背起十字架。这些容易吗？\par

\BibleRef{玛 14:31}。连伯多禄的信心也动摇，开始下沉。\par

\BibleRef{玛 15:19-20}。从心里发出恶念等等。\par

\BibleRef{玛 16:25}。凡要丧失自己生命的，必得着生命。\par

\BibleRef{玛 18:7}。\BibleQuote{那绊倒人的有祸了。} 但是\par

\BibleRef{雅 3:2}。我们在许多事上都有过犯。\par

\BibleRef{斐 2:21}。众人都求自己的事。\par

\BibleRef{玛 19:21}。那年轻的法学士遵守了全部法律，却还是失败了。\par

\BibleRef{玛 23:26}。对法利塞人的祸哉，按着他们的分量临到众人身上。\par

14.\par

\BibleRef{玛 26:39}。\Quote{不要照我的意愿，而要照你的意愿。}然而克里托布鲁斯说，凭他自己的意志，他能行义。\par

\BibleRef{谷 14:37}。\Quote{你们竟不能同我一起醒寤一个时辰吗？}他们\Emph{不能}。\par

\BibleRef{谷 6:5}。他不能行任何大能的事，因为他们不信。\par

\BibleRef{谷 7:24}。\Quote{他进入了提洛和漆冬的边界。}如果基督不能按他所愿的去做，我们怎能呢？\par

\BibleRef{谷 9:5}。伯多禄在显圣容时的请求显示了他的无知。\par

\BibleRef{谷 13:32}。连圣子也不知道一切事，我们怎能呢？\par

\BibleRef{谷 14:35}。如果可以的话。你怎么能说每时每刻都有可能避免犯罪呢？\par

15.\par

\BibleRef{谷 16:14}。连宗徒们也显出不信和心硬。\par

\BibleRef{若壹 5:19}。世界躺卧在那恶者里面。\par

\BibleRef{路 1:20}。连匝加利亚也不信天主的信息。\par

\BibleRef{玛 17:15}。门徒们不能治好那个癫痫病人，因为他们不信。\par

\BibleRef{谷 4:34}。门徒们关于地位的争论。\par

\BibleRef{路 9:54}。雅各伯和若望显出报复的心态。\par

\BibleRef{路 14:26-27}。舍弃一切和背十字架的命令并不容易。\par

\BibleRef{路 16:15}。在人前被高举的，在天主眼中却是可憎的。\par

\BibleRef{路 17:1}。绊倒人的事是免不了的。\par

\BibleRef{路 17:6}。宗徒们的信德连一粒芥菜子都不如。\par

\BibleRef{雅 3:2}。\par

\BibleRef{玛 17:19}。\par

16.\par

\BibleRef{路 18:1}。我们应该常常祈祷。这显示出我们的软弱。\par

\BibleRef{路 18:27}。那么，谁能得救呢？这是可能的，但唯独对天主是可能的。\par

\BibleRef{路 22:24}。最后晚餐时关于谁为大的争论。\par

\BibleRef{路 22:31-32}。伯多禄的信德几乎被撒殚胜过。\par

\BibleRef{路 22:43}。连基督在祂的忧苦中也需要一个天使来加祂力量。\par

\BibleRef{路 22:46}。祈祷，免得进入诱惑。\par

17.\par

\BibleRef{若 5:30}。连基督也说：\Quote{我什么也不能凭自己做什么}；而且\par

\BibleRef{若 7:10}。祂在上帐棚节的事上显得犹豫不决。\par

\BibleRef{若 7:19}。你们中没有人遵行法律。\par

\BibleRef{若 8:3}。控告那犯奸淫的妇人的众人，没有一个是没有罪的。基督在地上写了他们的名字\BibleRef{耶 17:13}。\par

\BibleRef{若 10:8}。凡在基督之前来的\BibleRef{耶 14:15}都是贼。\par

\BibleRef{若 17:12}。我保全了他们——他们并没有保全自己。\par

\BibleRef{宗 15:39}。保禄和巴尔纳伯争吵了。\par

\BibleRef{宗 16:6-7}。他们被禁止在他们选择的地方讲道。\par

\BibleRef{宗 16:18}。连宗徒们在全备的光照下，也显出他们对恩宠的依赖。\par

\BibleRef{宗 17:30}。基督以前的时期是无知的时期。\par

\BibleRef{格前 4:19}。我\Emph{如果主愿意}就到你们那里去。\par

\BibleRef{雅 2:10}。在一件事上跌倒，就是违反全部诫命而有罪。\par

\BibleRef{雅 3:2}。我们在许多事上全都跌倒。\par

\BibleRef{雅 3:8}。舌头是致命的毒物。\par

19.\par

\BibleRef{雅 4:1}。战争起源于我们的私欲。达味确实说过：\par

\BibleRef{咏 25:2}。\Quote{求你察验我，考验我，}等等。这种自信导致了他的跌倒。\par

\BibleRef{咏 50:1}。天主，求你怜悯我。\par

\BibleRef{咏 79:5}。\Quote{你以眼泪的饼喂养我们。}类似地\par

\BibleRef{咏 30:6-7}。我曾说：我永不动摇……你掩面不顾。\par

\BibleRef{咏 31:5}。我说，我要承认我的过犯。\par

\BibleRef{咏 37:5-6}。\Emph{祂}要使你的正义如光出现。\par

\BibleRef{咏 36:39}。义人的救恩\Emph{来自主}。\par

\BibleRef{咏 37:7}。我肉体毫无健康。\par

\BibleRef{罗 7:18}。在我肉体中不住着善。\par

\BibleRef{咏 37:8}。拉丁通行本：我的腰充满诡诈。\par

\BibleRef{咏 38:5}。他使我们的年日如手掌宽。\par

\BibleRef{咏 68:5}。我的罪愆不能向你隐藏。\par

\BibleRef{咏 76:2}。我的灵不肯受安慰。\par

\BibleRef{咏 76:10}。这是至高者右手的变化。\par

20.\par

\BibleRef{咏 88:2}。仁慈将被\Emph{建立}直到永远。\par

\BibleRef{咏 90:6}。谁能脱离那\Quote{在黑暗中行走的东西}？因为\par

\BibleRef{咏 10:2}。\Quote{恶人弯弓搭箭}——异端者的形象。\par

\BibleRef{咏 91:14}。栽种在主殿中的，必发旺。\par

\BibleRef{咏 103:8-10}。主满有怜悯。\par

\BibleRef{撒下 8:13-14}。达味以谦卑承认软弱来领受应许。\Quote{天主啊，这是人的律法吗？}\par

\BibleRef{撒下 16:10}。他在阿彼瑟的暴力和史米的咒骂下谦卑自己。\par

\BibleRef{撒下 17:14}。并且只靠天主挫败阿希托费耳的计谋而得救。\par

\BibleRef{列上 14:8}。是天主将王国赐给了雅洛贝罕。\par

21.\par

\BibleRef{列上 15:11}。阿撒虽是个好人，却有缺陷。\par

\BibleRef{列上 19:4}。厄里亚从依则贝耳面前逃跑。\par

\BibleRef{咏 117:6}。主是我的保护者。\par

\BibleRef{编下 17:3}。约沙法特兴旺，因为主与他同在。然而\par

\BibleRef{编下 19:2}。他因与阿哈布联合而受责备。\par

\BibleRef{编下 22:9}。阿哈齐雅因是义人约沙法特的后裔，得以葬在君王中。\par

\BibleRef{列下 18:3-4,7}。希则克雅行了大事，但只靠主的帮助。\par

\BibleRef{列下 18:14}。他将圣金给了亚述王。\par

\BibleRef{列下 18:22}。连犹大最好的君王也不完美。\par

\BibleRef{列下 20:1,5}。希则克雅在死亡临近时哭泣，并藉着特殊的怜悯得以康复。\par

\BibleRef{列下 20:13,17}。但他因接待巴比伦的使者而犯罪。\par

\BibleRef{编下 32:26}。他因心高气傲而跌倒。\par

\BibleRef{编下 34:2}。约史雅是个义人；然而\par

\BibleRef{编下 34:22-23}。他需要胡耳达的帮助；而且\par

\BibleRef{编下 35:22}。他因不听天主的警告而被杀；而且\par

\BibleRef{编下 35:23}。先知们也是软弱有罪的。\par

\BibleRef{哀 4:20}。耶肋米亚哀叹他的跌倒。\par

\BibleRef{户 20:10,12}。梅瑟因在默黎巴的罪而受罚。这就是\BibleRef{咏 140:6}的意义。拉丁通行本：\Quote{他们的审判官被吞没，与磐石结合，等等。}\par

\BibleRef{欧 2:19}。天主以慈悲宽恕以色列的不忠。\par

\BibleRef{欧 11:9}。\Quote{我必不进入城中。}只有圣者不与不虔之众联合。\par

\BibleRef{亚 6:13}。我们将正义变为苦艾。\par

\BibleRef{纳 1:14}。水手们承认天主兴起风暴是公义的。\par

\BibleRef{米 7:2}。虔敬的人从地上灭绝了，等等。\par

\BibleRef{米 6:8}。行公义、好怜悯、谦卑地与天主同行的命令，唯有藉着谦卑的信德才可能，因为\par

\BibleRef{咏 139:6}。\Quote{恶人在各处游行，}而且\par

\BibleRef{雅 4:6}。天主赐恩宠给谦卑的人。\par

24.\par

\BibleRef{哈 3:16}。愿朽坏进入我的骨头，只要我能安息，等等。\par

\BibleRef{匝 3:1}。耶叔亚被描绘为穿着污秽的衣服，并借天主的仁慈而得释放。\par

\Quote{但是\Person{约维尼安}的继承人却说：我完全无罪，我没有污秽的衣服，我受自己意志的支配，我比宗徒还大。宗徒做他所不愿意做的，不做他所愿意做的；但我是做我愿做的，不做我不愿做的；天国已为我预备好了，或者说我以自己道德的生活为自己预备了它。\Person{亚当}应受惩罚，其他那些自认与\Person{亚当}违命相似而自觉有罪的人也如此；惟独我和我的同伙无所畏惧。别人关在斗室内，从未见过妇女，因为他们可怜！不听我的话，备受欲望的折磨；成群妇女环绕我，我却毫无私欲偏情。因为我正应验了这句话：\InnerQuote{圣石在地上滚来滚去}；我对罪的吸引力无动于衷的原因，就是带着基督的胜利纪念品在我身体里的自由意志的力量。} 但是让我们听天主藉\Person{依撒意亚}之口所宣告的（\BibleRef{依 3:12}）：\Quote{我的子民啊，那称你们有福的，使你们走迷了路，毁坏了你们所行的路径。} 谁是颠覆天主子民的最大人物——是那依靠自由抉择能力、藐视造物主帮助、满足于随从己意的人，还是那畏惧按主的诫命细节受审判的人？对于这种人，天主说：\Quote{祸哉，你们在自己的眼中看为智慧，在自己的眼中看为通达的人。} \Person{依撒意亚}，若我们按希伯来文，哀叹说：\Quote{祸哉，我！因为我沉默不言，因为我是一个嘴唇不洁的人；我住在嘴唇不洁的人民中间，因为我的眼看见了万军的上主。} 他因自己功德及道德的生活得以享见天主，且意识到自己的罪恶，承认自己有嘴唇不洁。不是因为他曾说过任何违逆天主旨意的话，而是因为他或出于恐惧，或出于深感羞耻，竟默不作声，没有像先知所应如此自由地斥责人民的错误。我们这些罪人何时责备过犯过者？我们谄媚财富，为不义之财而接受罪人的情面？因为我们很难说自己对需要其帮助的人完全坦率直言。假设我们不做他们所做的事，假设我们远离一切形式的罪；但闭口不言真理当然也是罪。然而在\WorkTitle{七十贤士译本}中，我们看到的不是\Quote{因为我沉默不言}，而是\Quote{因为我被刺痛}，即因意识到罪恶而刺痛；于是先知的话得以应验：\Quote{我的生命因被荆棘刺中而变为苦难。} 他被罪的荆棘刺中；你却以美德的花朵装饰自己。\BibleRef{依 24:21} \BibleQuote{那时，月亮要羞惭，太阳要蒙羞，因为上主在高天惩罚天上的军队。} 这由另一段经文解释。\BibleRef{约 25:5} \BibleQuote{群星在他眼中也不洁净}，还有\BibleRef{约 4:18} \BibleQuote{他指责自己的天使为愚昧。} 月亮羞惭，太阳蒙羞，天空披上麻衣；而我们却要无畏而快乐地，好像完全没有罪似的，面对审判者的威严吗？那时，山岳要溶化，即一切因骄傲而被举起的人，以及天上的万军，无论是星辰还是天使的权势，当诸天被卷起如同书卷，他们所有的军队都要像叶子一样凋落？\par

此后的论证主要通过引用先知书的段落进行。\par

25.\par

\BibleRef{依 34:5} \BibleQuote{我的刀剑在天上已喝足了，它要降在厄东。} 地上的罪更是何等激起愤怒！\Place{厄东}意为血，血不能承受天国（\BibleRef{格前 15:50}）。\par

\BibleRef{依 45:9} \BibleQuote{祸哉，那与他的造物主争辩的人。}\par

\BibleRef{依 53:6} \BibleQuote{我们都如羊走迷了路。}\par

\BibleRef{则 16:14} \Place{耶路撒冷}在美貌上完美无缺；然而\par

\BibleRef{则 16:60-61} 她的救恩不是因功绩，而是因怜悯。\par

\BibleRef{鸿 1:3} 他虽洁净，却不会使你无辜。\par

\BibleRef{格前 15:9} 我不配——因为我曾迫害。\par

\BibleRef{则 20:43-44} 当\Place{耶路撒冷}被赦免时，她仍将记念自己的罪。\par

让我们惭愧地承认，这些话乃是已经获得赏报的人所说出的；我们这些地上的罪人，且仍在脆弱、会死的身体中，让我们采用天上圣人的话，他们甚至已经领受了不朽性及不死性。\BibleRef{则 32:17} \BibleQuote{你们还说：上主的道路不公平，可是你们的道路才是不公平。} 将本应归于我们意志的罪归咎于造物主的不义，并诽谤祂的正义，这是法利塞人的骄傲。属灵圣殿（即教会）的祭司\Person{匝多克}的子孙，不穿着他们的礼服出去到人民中，免得因与人的交往而失去圣洁，受到玷污。而你，身处稠人广众之中，且是一个普通人，难道你以为自己是洁净的吗？\par

26. 让我们匆匆浏览\Person{耶肋米亚}先知：\par

\BibleRef{耶 5:1-2} 是否有行公义的人，等等。\par

\BibleRef{耶 7:21-22} 天主弃绝祭献，因为敬拜者的邪恶生活。\par

\BibleRef{耶 13:23} \BibleQuote{古实人能改变他的皮肤吗？}\par

27.\par

\BibleRef{耶 17:14} \BibleQuote{上主，求你医治我。}否则\Person{耶肋米亚}只能如接着引用的经文所说，\par

\BibleRef{耶 20:14,17-18} 愿我出生的那日受咒诅，等等。\par

\BibleRef{耶 23:23} \BibleQuote{我岂是近处的天主，等等。}他是如此意识到天主的能力。\par

\BibleRef{耶 24:6-7} 天主，而非他们自己，将栽植他们，等等。\par

\BibleRef{耶 26:21-24} \Person{耶肋米亚}需要\Person{阿希甘}的帮助。何况我们更需要天主的帮助。\par

28.\par

\BibleRef{耶 31:34} 新盟约的许诺。\par

\BibleRef{耶 32:30} 以色列子民不断作恶。\par

\BibleRef{耶 37:18-19} 然而\Person{耶肋米亚}自己在\Person{漆德克雅}面前也战栗。\par

\BibleRef{耶 30:10-11} \BibleQuote{不要害怕，\Person{雅各伯}，因为\Emph{我与你同在}。}\par

29.\par

\BibleRef{亚 6:14} \BibleQuote{我们凭自己的力量为自己取了角。}这是异端者的夸口。但\par

\BibleRef{依 16:6} 他的力量（\Place{摩阿布}的力量）绝比不上他的骄傲。\par

\BibleRef{耶 1:7,20} 人的罪只有因天主施恩才能消除。如果你放弃对自然能力的断言，我就承认你全部的论点成立，但仅靠天主的恩赐。\par

\BibleRef{哀 3:26} 人最好安静等候主的救恩。\par

30.\par

\BibleRef{达 4:17} 至高者在人的国中掌权。\par

\BibleRef{咏 113:7-8} 他从灰尘里抬举贫寒人。\par

\BibleRef{依 40:17} 他在天上和地上随意行事。\par

\BibleRef{加下 5:17} 的话说，\Person{安提约古·厄丕法乃}有能力推翻圣殿，\Quote{因为罪恶众多}，这段引文与达尼尔的忏悔相关联。\par

\BibleRef{达 9:5} \Quote{我们犯罪作恶}，这由……所显示。\par

20. \BibleRef{达 9:20} \Quote{我正祈祷的时候}等等，表明是个人的，不只是国家的忏悔。\par

24. \BibleRef{达 9:24} 七十周的预言表明先知只仰望天主来建立正义。\par

因此，直到那终结来临，这必朽坏和必死的要穿上不朽坏和不死，我们必然难免犯罪；并非如你错误地所说，由于我们本性和受造的缺陷，而是由于人意志的脆弱和善变，片刻之间就会变化；因为唯有天主不变。你问\Person{亚伯尔}、\Person{哈诺客}、\Person{农的儿子若苏厄}、\Person{厄里叟}以及其他圣人在哪些方面犯了罪。无需在灯心草中找结；我坦承我不知道；我只希望，当罪过明显时，我仍能保持沉默。\BibleRef{格前 4:4} \Quote{我虽然自觉良心无愧，}圣保禄说，\Quote{却不能因此就成义。} \BibleRef{撒上 16:7} \Quote{人看外貌，主却看人心。}在他面前无人能成义。所以保禄确信地说，\BibleRef{罗 3:23} \Quote{众人都犯了罪，亏缺了天主的荣耀；}以及\BibleRef{迦 3:22} \Quote{天主把众人都圈在罪里，为要怜悯众人。}其他经文也多次重复。\par

\SourceRecord{Translated by W.H. Fremantle, G. Lewis and W.G. Martley.；Nicene and Post-Nicene Fathers, Second Series；Vol. 6.；Edited by Philip Schaff and Henry Wace.；Buffalo, NY: Christian Literature Publishing Co.,；1893.；<http://www.newadvent.org/fathers/30112.htm>.}

% structure: p0001 source=h1 target=section reason=page-title

\section{驳斥佩拉纠派（卷三）}

1. \Person{Critob}. 我陶醉于您口才的丰沛，但同时我要提醒您，\BibleRef{箴 10:19}\BibleQuote{多言多语难免无过。} 这和我们眼前的问题有什么关系？您肯定承认那些领受基督徒圣洗的人是没有罪的。既然没有罪，他们就是正义的。而一旦正义，他们若能小心，就能保持正义，从而一生避免一切罪。\par

\Person{Attic}. 您不因追随已被揭穿和谴责的\Person{约维尼安}的意见而羞愧吗？因为他所依据的正是和您同样的证据和论证；不，更确切地说，您热切追求他的发明，想要在东方宣讲那些曾在罗马、不久前又在非洲被谴责的论调。那么，请阅读对他的答复，您会在那里找到对您自己的答复。因为在教义和争议点的讨论中，我们必须关注的是事而不是人。但让我告诉您，圣洗宽恕过去的过犯，并不能保持未来的正义；保持正义取决于努力和勤奋，以及热忱，尤其是天主的仁慈。祈求是我们的本分，赐予我们所求的是祂的本分；开始是我们的，完成是祂的；献上我们所能的是我们的，成就我们所不能的是祂的。\BibleQuote{若不是上主建造房屋，建造的人也是徒然劳苦；若不是上主守护城池，守护的人也是徒然警醒。} 因此，宗徒\Person{保禄}\BibleRef{格前 9:24}吩咐我们要这样跑，好能得到奖赏。固然都跑，但只有一个得奖赏。在圣咏中记载：\BibleQuote{上主，你以你的恩惠像盾牌一样给我们加冕。} 因为我们的胜利是藉祂的保护和祂的盾牌赢得的，胜利的冠冕也是这样获得的；我们在此奔跑，好能在将来得到；谁在这世上成为得胜者，谁就在那里领受冠冕。当我们受洗后，我们被告知：\BibleRef{若 5:14}\BibleQuote{看，你已经痊愈了，不要再犯罪，免得你遭遇更不幸的事。} 又说：\BibleRef{格前 3:16}\BibleQuote{你们不知道你们是天主的宫殿，天主圣神住在你们内吗？谁若亵渎天主的宫殿，天主必要毁灭他。} 在另一处：\BibleRef{编下 15:2}\BibleQuote{上主与你们同在，你们若与他同在，他也要与你们同在；你们若离弃他，他也要离弃你们。} 您认为，有哪一个人，在他身上如同在圣所和至圣所中一样，基督的纯洁是恒久的，而其圣殿的宁静从未被罪的乌云所笼罩？我们不能总是有同样的面容，尽管\Person{苏格拉底}被哲学家们虚假地夸口说曾有那样的经历；我们的心灵怎能总是如一呢？正如人有多种面部表情，他们内心的情感也各不相同。如果我们能够一直浸在洗礼的水中，罪就会从我们头上飞过，不沾我们身。圣神会保护我们。但仇敌攻击我们，即便被战胜也不离开，而是不断埋伏，为要暗中射击心里正直的人。\par

2. 在希伯来人福音中——这福音是用加色丁语和叙利亚语写成，但使用希伯来字母，且纳匝肋人至今仍在使用（我指的是宗徒们的福音，或普遍所认为的玛窦福音，有一份抄本在\Place{凯撒勒雅}的图书馆里）——我们发现：\Quote{请看，我们主的母亲和祂的兄弟对祂说：\InnerQuote{\Person{洗者若翰}为赦罪而施洗，我们去吧，也由他受洗。} 祂对他们说：\InnerQuote{我犯了什么罪，竟要去由他受洗？除非我所说的这些话不过是出于无知。}} 在同一卷书中：\Quote{如果你的兄弟在言语上得罪了你，并向你赔补，你就一天内要原谅他七次。} 祂的门徒\Person{西满}对祂说：\Quote{一天七次？} 主回答他说：\Quote{我对你说：直到七十个七次。} 即使是先知们在被圣神傅油之后，也曾有言语上的罪过。宗徒时代的殉道者\Person{依纳爵}大胆地写道：\Quote{主拣选了宗徒，他们比众人更是罪人。} 正是关于他们的迅速悔改，圣咏作者唱道：\BibleQuote{他们的软弱增多了，然后就急忙转向。} 如果你不承认这证据的权威，至少承认它的古老，并看看所有好教会人士的意见。假设一个受过洗的人，或在受洗时，或在受洗当天就去世了，我愿意大方地承认，他既没有想也没有说任何导致他因错误和无知而陷入罪过的话。那么，难道他就因此而没有罪，因为他似乎不是胜过了罪，而是避免了罪？真正的理由难道不是因天主的仁慈，他从罪的牢狱中被释放，并归向了主吗？我们也这样说：天主能行祂所愿的；而人靠自己、凭自己的意志，不能如你所主张的那样没有罪。如果他能，那么你现在提到恩宠一词就是徒劳，因为有了这样的能力，他就不需要恩宠。但如果他不能在没有天主恩宠的情况下避免犯罪，那么你赋予他一种他所没有的能力，就是愚蠢的。因为凡事取决于他人意志的，就不在于你断言他有能力的那人，而在于显然不可或缺的帮助者。\par

\Person{C.} 你这乖僻，或者说无谓的争论，是什么意思？你难道连这一点都不肯承认——一个人离开洗礼的水时，是脱离罪恶的吗？\par

\Person{A.} 要么是我没有表达清楚我的意思，要么是你理解迟钝。\par

\Person{C.} 怎会如此？\par

A. 请记住您所主张的以及我所说的。您争辩说，人若愿意，就能无罪。我回答说，这是不可能的；这并不是说我们认为人在圣洗后不是立即无罪，而是说那段无罪的时间决不可归功于人的能力，而应归功于天主的恩宠。所以，不要将能力归于人，我就承认事实。因为一个靠自身不能的人，怎么能有能力呢？或者，那取决于身体立即死亡的无罪是什么呢？倘若人的生命延长，他必定会陷入罪恶和无知。\par

C. 您的逻辑堵住了我的嘴。您不是以基督徒的纯朴说话，而是用一些关于\Quote{存在}和\Quote{能够存在}的精微区别来纠缠我。\par

A. 难道是我在玩弄这些文字把戏吗？那篇文章出自您自己的作坊。因为您说，人不是无罪，而是能够无罪；而我却要承认您所否认的，即人因天主的恩宠而无罪，同时坚持他靠自己不能做到。\par

C. 如果诫命我们不能遵守，那么颁布诫命就是无用的。\par

A. 没有人怀疑天主命令了可能的事。但因为人不做他们所能做的，所以全世界都处于天主的审判之下，并需要祂的怜悯。另一方面，如果您能提出一个遵守了全部法律的人，那么您当然能证明有一个人不需要天主的怜悯。因为一切可能发生的事，要么发生在过去、现在，要么将来。至于您声称人若愿意就能无罪，请证明这在过去发生过，或至少现在确实发生；未来自会显明。然而，如果您不能指出任何一个完全无罪的人（无论是现在还是过去），那么我们就只能将讨论限定在未来。同时，在三个时间段中的两个（过去和现在）里，您已失败被俘。如果将来有谁比先祖、先知、宗徒更伟大，因为他没有罪，那么您或许能说服未来世代的人关于他们的时代的事。\par

4. C. 随您怎么说，随您怎么辩，您永远不能从我这里夺走天主一次赐予的自由意志，也不能剥夺天主所赐的、只要我愿意就能的能力。\par

A. 让我们举一个证据为例：\BibleQuote{我找到了叶瑟的儿子达味，一个合我心意的人，他要承行我的一切旨意。} 毫无疑问，达味是一个圣洁的人，但那位被拣选来承行天主一切旨意的人，却因某些行为而受责备。当然，那被拣选来承行天主一切旨意的人是有可能做到的。天主也没有错，祂预先说他会承行祂的一切旨意如同命令，但错在于他没有做所预言的事。因为天主并没有说祂找到了一个会不折不扣地遵行祂的命令、实现祂旨意的人，而只说他会承行祂的一切旨意。我们同样说，人如果愿意，可以避免犯罪，这取决于他当地、当时的环境和身体的软弱，只要他的心志专注于正义，琴弦在竖琴上绷紧。但如果人稍一松懈，就像逆水划船的船夫，一旦稍一停歇，船就会向后滑，被流水带到他不愿去的地方。人的情形就是这样；如果我们稍一疏忽，就认识到自己的软弱，发现自己的能力是有限的。您以为宗徒保禄在写 \BibleRef{弟后 4:13} \BibleQuote{我在特洛阿留在卡尔颇那里的外衣（或斗篷），你来时带来，还有那些书，尤其是那些羊皮卷} 时，是在思念天上的奥秘，而不是那些日常生活所需、满足身体需要的东西吗？请给我找一个从不饥饿、口渴、寒冷，不知疼痛、发烧或小便刺痛之苦的人，我就承认人能只想着德行。当宗徒被仆人打时，他这样回答那下令打人的大司祭：\BibleQuote{天主必要打击你，你这粉刷的墙。} 我们怀念救主的忍耐，祂像牵到宰杀之地的羔羊，不开口，却仁慈地对打祂的人说（\BibleRef{若 18:23}）\BibleQuote{我若说得不对，你指证不对的地方；若说得对，你为什么打我？} 我们并非贬低宗徒，而是显扬天主的荣耀，祂在肉身内受苦，并克服了加在肉身和肉身软弱上的恶——更不用说宗徒在别处所说的话（\BibleRef{弟后 4:14}）\BibleQuote{铜匠亚历山大加给了我许多恶事；主，公义的审判者，要在那一日报应他。}\par

5. C. 我一直渴望说些什么，但话到嘴边又止住了。您逼我说出来。\par

A. 谁阻止您说出您的想法？要么您要说的是好的——您不应该剥夺我们好的东西；要么是坏的，那么，您保持沉默不是出于顾及我们，而是因为羞愧。\par

C. 我要说，我终究要说出我的想法。您的整个论证都指向一点：您指控本性，并责怪天主把人造成这样。\par

A. 这就是您想说又不想说的话吗？请说出来，好让大家都能从您的智慧中受益。您是在指责天主把人造成人吗？也让天使抱怨他们是天使吧。让每个受造物讨论：它为什么被造成它被造的样子？而不是讨论造物主可能把它造成什么。我现在必须用童年时的修辞练习来自娱，从蚊虫和蚂蚁到革鲁宾和色辣芬，探究为什么每个受造物没有被赐予更幸福的命运。当我到达那崇高的权能时，我要争论这一点：为什么只有天主是唯一的天主，而没有使万物都成为神？因为按照您的说法，祂要么不能这样做，要么就是嫉妒。指责祂，质问祂为什么允许魔鬼存于这个世界，并在您赢得胜利时夺走花冠。\par

C. 我不至于愚蠢到抱怨魔鬼的存在，因为藉着他的恶意，死亡进入了世界；但令我痛心的是：教会的显贵们，以及那些僭取导师称号的人，毁灭了自由意志；一旦自由意志被毁灭，就为摩尼教徒打开了道路。\par

A. 难道我就是自由意志的毁灭者，因为在整场讨论中，我唯一的目标就是既维护天主的全能，又维护自由意志吗？\par

C. 你怎么可能拥有自由意志，却又说人没有天主的助佑就什么也不能做呢？\par

A. 如果谁将自由意志与天主的助佑相连，就要受责备，那么我们就应该赞美那些废除了天主助佑的人。\par

C. 我并非使天主的助佑成为多余，因为我们的全部能力都来自祂的恩宠；但我和那些与我见解相同的人，将两者保持在各自的界限内。我们将自由抉择的能力之恩赐归于天主的恩宠；将行事或不行事归于我们自己的意志；这样，行善或不行善的赏罚就可以得到维持。\par

6. A. 在我看来，你似乎陷入了遗忘，并且正在重复已经讨论过的论点，仿佛之前什么话都没有说过。因为，通过这场长久的讨论，已经确定：主以祂赐予我们自由抉择的同一恩宠，在我们个别的行动中协助并支持我们。\par

C. 那么，祂为什么要加冕并赞美祂自己在我们内所成就的事呢？\par

A. 那就是说，祂赞美我们的意志，因为它奉献了它所能的一切；赞美我们的辛劳，因为它努力行动；赞美我们的谦卑，因为它始终仰望天主的助佑。\par

C. 那么，如果我们没有做祂所命令的事，要么天主愿意协助我们，要么祂不愿意。如果祂愿意并协助了我们，而我们却没有做我们所愿意的，那么被战胜的就是祂，而不是我们。但如果祂不愿意帮助，那么那个愿意遵行祂旨意的人就不应受责备，而应责备那能帮助却不愿意的天主。\par

A. 难道你没有看出你的两难推理使你陷入了亵渎的深渊吗？无论你取哪一边，天主要么是软弱的，要么是恶意的；这样，祂与其因是善的作者并给予帮助而受赞美，不如因没有阻止邪恶而受辱骂。那么，就责备祂吧，因为祂允许魔鬼存在，并且容忍了，并且仍然容忍世界上发生邪恶。这正是马西昂所问的，以及那整群割裂\WorkTitle{旧约}的异端者所问的；他们大多编造了类似这样的论证：要么天主知道人被安置在乐园里会违犯祂的命令，要么祂不知道。如果祂知道，那么人就不应受责备，因为他无法避免天主的预知；应受责备的是那创造了他、使他无法逃脱天主的预知的那一位。如果祂不知道，那么剥夺祂的预知，你就同时剥夺了祂的神性。根据同样的说法，天主将因拣选了撒乌耳而应受责备，因为他后来成了最坏的君王之一。救主也必须被定为无知或不义，因为祂在福音中说：\BibleRef{若 6:70}\BibleQuote{我岂不是拣选了你们十二人，而你们中间有一个是魔鬼吗？}请问祂为什么拣选犹达斯这个叛徒？为什么祂知道他是贼，却把钱袋交给他？要我告诉你原因吗？天主判断的是现在，而不是未来。祂不使用祂的预知来定一个人的罪，虽然祂知道此人将来会使祂不悦；但祂的良善和说不尽的仁慈是这样：祂拣选一个祂看出暂时会是好人、而祂知道将来会变坏的人，这样给他悔改和皈依的机会。这就是宗徒的意思，他说：\BibleRef{罗 2:4}\BibleQuote{难道你不知道天主的仁慈引领你悔改吗？但你凭着你顽固和不悔改的心，在忿怒的日子、天主公义审判显示的日子，为自己积蓄忿怒；祂要照每人的行为报应各人。}因为亚当犯罪，并不是因为天主知道他会这样做；但天主，作为天主，预知了亚当会凭自己的自由抉择做什么。你也可以指责天主撒谎，因为祂藉约纳的口说：\BibleQuote{再过三天，尼尼微就要毁灭了。}但天主会藉耶肋米亚的口回答：\BibleRef{耶 18:7}\BibleQuote{我何时对一个民族或一个王国发言，说要拔除、拆毁、毁灭它；如果我所针对的那个民族转离他们的邪恶，我就后悔我本想加给他们的灾祸。我何时对一个民族或一个王国发言，说要建立、栽植它；如果它在我眼中行恶，不听从我的声音，我就后悔我本想施予它们的恩惠。}有一次，约纳愤慨，因为按照天主的命令，他似乎说了假话；但他的忧愁被证明是毫无根据的，因为他宁愿说真话而使无数人灭亡，还是宁愿说假话而使他们得救？他的处境被这样阐明：\BibleRef{纳 4:10}\BibleQuote{你为那棵蓖麻（或葫芦）忧伤，它你没有劳苦，也没有使它生长，一夜发生，一夜死去；而我不该怜惜尼尼微这座大城，其中有十二万多人不能分辨左右手吗？}既然有如此多的孩童和纯朴的人（你永远无法证明他们是罪人），那么关于那些处于不同年龄阶段的男女居民，我们该说什么呢？根据斐罗，以及最聪明的哲学家柏拉图（正如\WorkTitle{蒂迈欧篇}告诉我们的），从婴儿期到衰老期，我们经历七个阶段；这些阶段如此逐渐而温和地相互跟随，以至于我们完全感觉不到变化。\par

C. 你整个论证的要点是这样的：希腊人称之为\OriginalTerm{自由意志}{\GreekText{αὐτέξουσιον}}、我们称之为自由意志的东西，你口头上承认，但实际上却毁掉了。因为你使天主成为罪的作者，声称人凭自己什么也不能做，而必须得到天主的帮助，我们所作的一切都归咎于祂。但我们的说法是：无论人行善或行恶，都因自由抉择的能力而归咎于他，因为他做了他所选择的，而不是归咎于那一次性地赐予他自由抉择的那一位。\par

A. 你的含糊其辞毫无用处；你已陷入真理的罗网。因为照此说法，即使祂自己不施予助佑，按照你的观点，祂也是邪恶的创造者，因为祂本可阻止却未阻止。有一句古老的格言：如果一个人能救他人免于死亡却没有救，他就是杀人者。\par

C. 我撤回并让步；你赢了；不过，这种胜利是用花言巧语颠覆真理，也就是说，不是靠真理，而是靠谎言。因为我能以宗徒的话回答你：\BibleRef{格前 11:6} \BibleQuote{我拙于言辞，却不拙于知识。}当你讲话时，你的修辞技巧令我难以招架，我似乎同意你；但当你停止讲话时，这一切便从我脑中消失，我清楚地看到，你的论证并非源于真理和基督徒纯朴的泉源，而是建立在哲学家们绞尽脑汁的诡辩之上。\par

A. 那么，你希望我再次诉诸圣经的见证吗？如果这样，你门徒们的夸口——说无人能反驳你的论证或解决你提出的问题——又作何解释？\par

C. 我不但希望，而且渴望你这样做。请给我指出圣经中任何一处，表明人丧失了自由选择的能力后，能凭自己去做他原本不愿或不能做的事。\par

8. A. 我们必须使用圣经的话，并非如你所提议，而是按照真理和理性的要求。雅各伯在祈祷中说：\BibleQuote{若上主天主要与我同在，在我所行的路上保护我，赐我食物吃，衣服穿，使我平安回到我父亲的家，上主就必作我的天主；我所立为柱子的这块石头，必成为天主的殿；凡祢所赐给我的，我必献给祢十分之一。}他没有说：若祢保存我的自由选择，我凭劳苦获得食物和衣服，回到我父亲的家。他将一切都交托于天主的旨意，以便配得他所祈求的。雅各伯从美索不达米亚回来时，\BibleRef{创 32:2} 天主的军旅——被称为天主的营——迎接他。他后来以一个男人形像的天使角力，并被天主坚固；于是，他不再被称为\Emph{欺骗者}雅各伯，而是得名\Emph{天主的至正直者}。因为他若未蒙主的助佑坚固和保障，就不敢回到他残忍的哥哥那里。接着我们读到：\BibleQuote{当他过了培奴耳后，太阳升起照在他身上，}培奴耳被解释为\Emph{天主的面容}。因此梅瑟也说：\BibleQuote{我面对面见了上主，我的性命仍得保全，}这不是出于本性的特质——而是出于天主的俯就，祂仁慈宽恕。这样，当天主使祂的面容光照我们，赐予我们力量时，正义的太阳便在我们身上升起。若瑟在埃及被囚于监牢，后来我们得知狱吏信任他的忠诚，将一切交在他手中。原因如下：\BibleRef{创 39:23} \BibleQuote{因为上主与他同在，凡他所做的，上主都使之顺利。}因此，也给法老的臣仆启示了梦境，法老做了一个无人能解的梦，好让若瑟被释放，他的父亲和弟兄得饱食，埃及在饥荒中得救。此外，天主在夜间异象中对以色列说：\BibleRef{创 46:3} \BibleQuote{我是你父亲的天主，不要害怕下到埃及去，因为我要使你在那里成为一大民族；我要同你一起下到埃及，也必定带你上来；若瑟要亲手合上你的眼。}在这段经文中，自由选择的能力在哪里？他敢于到儿子那里去，将自己托付给一个不认识上主的民族，这岂不完全是靠他父亲的天主的助佑？百姓以强有力的手和伸展的臂膀从埃及被释放；那不是梅瑟和亚郎的手，而是那位藉着神迹奇事释放百姓，最后击杀埃及长子，使那些最初顽固扣留百姓的人急切催促他们离开的天主的手。撒罗满说：\BibleRef{箴 3:5} \BibleQuote{你要全心信赖上主，不要依赖自己的聪明；应步步体认祂，祂必修平你的行径。}请明白祂所说的——我们不可倚靠自己的智慧，只应倚靠上主，因人的脚步由祂引导。最后，我们奉命向祂显示我们的道路，并让祂知晓，因为我们的道路不是靠自己的辛劳修直，而是靠祂的助佑和仁慈。因此经上记着：\BibleQuote{求祢在我面前修直我的道路，}这样祢所视为正直的，在我眼中也视为正直。撒罗满又说同样的话：\BibleRef{箴 16:3} \BibleQuote{将你的作为交托于上主，你的计划必能成功。}当我们把我们一切所为都交托于上主——我们的助佑者，如同建立在坚固的磐石上，并将一切归功于祂时，我们的计划就成功了。\par

9. 宗徒保禄迅速数算天主的恩惠，以\Quote{这些事谁能胜任？}结束。\BibleRef{格后 2:16} 因此，他在另一处\BibleRef{格后 3:4} 也说：\Quote{我们藉着基督在天主面前才有这样的信心。这不是说我们凭自己能够承担什么事，好似出于自己；我们之所以能够承担，乃是出于天主，祂使我们能够成为新约的仆役，不是属于文字，而是属于圣神，因为文字叫人死，圣神却叫人活。}\newline{}我们还能以自由意志自夸，并滥用天主的恩惠来羞辱恩赐者吗？而同一蒙选的器皿明明\BibleRef{格后 4:7} 写道：\Quote{我们把这宝贝放在瓦器里，是为显明这卓然的力量是属于天主，而非出于我们。}因此，在另一处，他制止异端者的狂妄，\BibleRef{格后 10:17} 说：\Quote{夸耀的，应当因主夸耀。因为不是自我举荐的是可取的，而是主所举荐的。}又\BibleRef{格后 12:11} 说：\Quote{我虽然在一切事上都不落后于那些超级宗徒，但我什么也不是。}\Person{伯多禄}因所见奇事的伟大而惊慌，对主说：\BibleRef{路 5:8}\Quote{主，请离开我，因为我是个罪人。}主对他的门徒说：\BibleRef{若 15:5}\Quote{我是葡萄树，你们是枝条；那住在我内，我也住在他内的，他就结许多果实，因为离了我，你们什么也不能做。}正如葡萄树的枝条一旦与主干分离，立刻枯萎，同样，人的一切力量若失去天主的帮助，便衰败消灭。\Quote{没有人，} \BibleRef{若 6:44} 祂说，\Quote{能到我这里来，除非那派遣我的圣父吸引他。}当祂说\Quote{没有人能到我这里来}时，便击碎了自由意志的骄傲；因为即使人愿意走向基督，除非那后续的——\Quote{除非我天上的圣父吸引他}——得以实现，否则愿望徒然，努力枉费。同时应注意，被吸引的人并非自由奔跑，而是被带领，或因他退缩迟延，或因他不愿前行。\par

10. 现在，一个不能凭自己的力量和劳苦来到耶稣面前的人，怎能同时避免一切罪过？并且永久地避免，为自己求取属于天主大能的称号呢？因为如果祂和我都无罪，我和天主之间有什么区别？我只再引证一个证据，以免使你和听众感到厌倦。\BibleRef{艾 6:1} 睡眠从\Person{阿哈随鲁}（七十贤士本称为\Person{阿尔塔薛西斯}）的眼中被挪去，使他翻查忠臣的备忘录，发现\Person{摩尔德开}（此人曾揭发一桩阴谋而救他一命），从而使得\Person{艾斯德尔}更蒙悦纳，整个犹太民族得以逃脱迫近的死亡。无疑，那位统辖整个东方、从印度直到北方和埃塞俄比亚的威严君王，在享用来自世界各地珍馐佳肴之后，理当想要睡觉、休息，顺从自己对睡眠的自由选择，若不是万善的供应者天主阻止了自然的进程，以致违背自然，暴君的残酷才得以被克服。若我试图列举圣经中的所有实例，就太冗长了。圣人们所说的一切都是向天主的祈祷；他们全部的祈祷和恳求都是对天主怜悯的奋力争取，好让我们这些凭自身力量和热忱无法得救的人，能因祂的仁慈得以保存。然而，当我们涉及恩宠和怜悯时，自由意志部分失效；我说部分，因为在于它的是我们愿意、渴望并赞同我们所选择的道路。但在于天主的是我们是否有能力靠祂的力量和帮助来实现我们所渴望的，并完成我们的辛劳和努力。\par

11. C. 我只是说，我们在个别行动中，是在创造和法律的恩宠中找到天主的帮助，而非在每个行动中，以免自由意志被摧毁。但我们中有许多人坚持说，我们所做的一切都是在天主的帮助下完成的。\par

A. 持这种说法的人必须离开你们的团体。那么，要么你自己这样说而加入我们一边，要么你如果拒绝，你就和那些不同意我们观点的人一样成为我们的敌人。\par

C. 如果你说出我的感想，我就站在你一边；或者更确切地说，如果你不反对我的感想，你就站在我一边。你承认身体的健康，却否认灵魂的健康（灵魂比身体更强）。因为罪之于灵魂，如同疾病或创伤之于身体。那么，如果你承认就肉体而言，人可能是健康的，为何不说就精神而言，他可能是健康的呢？\par

A. 我要沿着你所指出的路线前进，\par

而你今日\newline{}绝不能逃脱；无论你呼唤到哪里，我都来。\par

C. 我洗耳恭听。\par

A. 我这是在向聋者说话。因此，我要回应你的论证。由灵魂与身体所组成的人，兼具两种实体的本性。正如身体若不为任何软弱所困扰，便可称为健康；同样，灵魂若不动摇、不烦乱，便没有过错。然而，虽然身体可能健康、健全、活跃，所有官能都充满活力，但它仍会时不时地遭受各种病痛的折磨；无论它多么强壮，有时也会因各种体液而感到不适；同样，灵魂承载着思绪与激动的冲击，即便避免船毁人亡，也并非无危险地航行；它记着自己的软弱，常为死亡而忧虑，正如经上所载：\Quote{什么人能生活而免见死亡？}——这死亡威胁所有必死之人，不是因本性的衰败，而是因罪恶之死，按先知的话说：\BibleRef{则 18:4}\Quote{犯罪的人，他必死亡。}此外，我们知道哈诺客和厄里亚尚未见过这种人与畜生共有的死亡。请给我指出一个从不生病、或病后永远健康平安的身体，我便给你指出一个从未犯罪、且在修德后永不再犯罪的灵魂。这是不可能的，尤其当我们想到恶习与美德毗邻，稍一偏离便会走入歧途或坠入深渊时，更是如此。顽固与坚忍、吝啬与节俭、慷慨与挥霍、智慧与狡猾、无畏与鲁莽、谨慎与胆小之间是多么微小的差距！其中一些被列为善，另一些被列为恶。身体的情况也是如此。你若防胆汁过多，痰液就会增加；你若过于快速地使体液干涸，血液就会发热并因胆汁而败坏，脸色便显蜡黄。毫无疑问，无论我们多么尽力遵循医生的所有建议，调节饮食，避免消化不良和任何滋生疾病的因素——这些疾病的根源有些为我们所隐藏，只有天主知道——我们仍会因寒冷而颤抖，因发热而燃烧，因腹痛而嚎叫，并祈求真正医生、我们的救主的帮助，并与宗徒们一起说：\BibleRef{玛 8:25}\Quote{主啊！救我们，我们丧亡了。}\par

C. 就算在童年、青年和壮年都无法避免一切罪恶；你能否认有许多义人和圣人，在陷入恶习之后，已全心投入修德并借此逃脱了罪恶吗？\par

A. 这正是我一开始对你说的——犯罪或不犯罪，行善或作恶，都取决于我们自己；因此，自由意志得以保存，但须视情况、时间及人性的软弱状态而定；然而，我们坚持认为，永无罪恶只保留给天主，以及那作为圣言降生成人、却没有沾染肉身的缺陷与罪恶的那一位。因为我能在短时间内避免罪恶，你不能逻辑地推断我能持续如此。我能不断地禁食、守夜、行走、歌唱、坐着、睡觉吗？\par

C. 那么，为什么圣经激励我们追求完美的义德呢？例如：\BibleRef{玛 5:8}\Quote{心里洁净的人是有福的，因为他们要看见天主。}以及\Quote{行为成全、遵行上主法律的人，是有福的。}天主对亚巴郎说：\BibleRef{创 17:1}\Quote{我是你的天主，你当在我面前行走，作个成全的人，我要在我与你之间立我的约，并要大大地增加你。}如果圣经所见证的是不可能的事，那么命令人去实行就是徒劳的。\par

甲：你玩弄圣经，直到把问题问得破旧不堪，使我想起变戏法的人在台上的伎俩：他扮演各种角色，一会儿是战神，一会儿是爱神；那起初全副严厉凶暴的人，化作了女性的温柔。因为你以貌似新颖的姿态提出的异议——\Quote{心里洁净的人是有福的}，\Quote{在道路上无玷的人是有福的}，以及\Quote{你要无瑕无疵}等等——被宗徒的回答所驳斥：\BibleRef{格前 13:9}\BibleQuote{我们现在所知道的只是部分，我们作先知所讲的也只是部分}，以及\BibleQuote{我们现在是藉着镜子观看，模糊不清，但到那完全的来到，那部分的就要被废除。}因此我们只有洁净之心的影子和肖像，这心将来注定要看见天主，并脱离一切斑点瑕疵，与亚巴郎一同生活。无论族长、先知或宗徒多么伟大，我们的主和救主都对他们说：\BibleRef{玛 7:11}\BibleQuote{你们纵使不好，尚且知道把好东西给你们的儿女，何况你们在天之父，岂不更将好东西赐给求祂的人？}然后，就连亚巴郎——曾听见对他说\Quote{你要无瑕无疵}——也因意识到自己的软弱，俯伏在地。当天主对他说：\Quote{你的妻子撒辣依不再叫撒辣依，要叫撒辣；我要由她赐给你一个儿子，我要祝福他，使他成为一大民族，人民的君王要由他而生}之后，叙述紧接着说：\Quote{亚巴郎遂俯伏在地，心里笑着说：\InnerQuote{一百岁的人还能生子吗？撒辣已九十岁，还能生子吗？}}亚巴郎对天主说：\Quote{但愿依市玛耳在你面前生存！}天主说：\Quote{不，你的妻子撒辣要给你生个儿子，你要给他起名叫依撒格。}等等。他确曾听见天主的话：\Quote{我是你的天主，你要在我眼中蒙悦纳，并且无瑕}；那么为什么他不相信天主所应许的，为什么他心里发笑，以为能瞒过天主，不敢公开笑出来呢？而且他提出不信的理由说：\Quote{一百岁的人怎能从九十岁的妻子得子呢？}\Quote{但愿依市玛耳在你面前生存！}他说。\Quote{依市玛耳是你曾经赐给我的。我不求难事，我满足于已得的祝福。}天主用一个奥秘的回答说服了他。祂说：\Quote{是。}意思是：你以为不会发生的事，必将发生。你的妻子撒辣要给你生个儿子，在她尚未怀孕、孩子尚未出生之前，我要给这孩子起名。因为，由于你暗中发笑的错误，你的儿子要叫依撒格，即\Emph{笑声}。但你若认为天主在今世被心里洁净的人看见，那么为何梅瑟——他曾说\Quote{我见了上主的面，我的生命得以保存}——后来又恳求能清楚地看见祂呢？而且因为他曾说他看见了天主，上主告诉他：\BibleRef{出 33:20}\BibleQuote{你不能看见我的面，因为人见了我，就不能再活了。}因此宗徒也称祂为那唯一看不见的天主，住在不可接近的光中，人未曾看见，也不能看见。而若望福音作者以神圣的声音作证说：\Quote{从来没有人见过天主；只有在父怀里的独生子，祂详述了祂。}那看见的，也详述了，不是所见者有多大，也不是详述者知道多少，而是按凡人所能领受的程度。\par

13. 而你认为在道路上无玷、遵行祂法律的人是有福的，你当用后一句解释前一句。从我引证的许多证据，你已经知道没有人能满全法律。如果宗徒为了赢得基督，与基督的恩宠相比，将从前在法律下看为盈利的视为粪土，那么我们就更应确知，基督和福音的恩宠之所以被加上，是因为在法律下没有人能成义。如今，若在法律下没有人成义，那么那仍在道上行走、奔向终点的人，怎能是在道路上完全无玷呢？的确，那在赛程中、在路上前进的人，不如那已抵达终点的人。那么，若仍在道上行走、在法律中前进的人已是无玷和完全的，那么那已抵达生命和法律终点的人，将更拥有什么呢？因此，宗徒论及我们的主，说到在世界终结、诸德圆满之时，祂将把圣洁的教会呈献给自己，没有瑕疵或皱纹；而你们却以为那仍在肉体——那属于死亡和朽坏的肉体——中的教会已是完全的。你们应当同格林多人一起被提醒：\BibleRef{格前 4:8}\BibleQuote{你们已经完美了，已经富足了，你们没有我们，自己就为王了；但愿你们真为王，好使我们也能与你们一同为王！}——因为真实无瑕的完全属于天上的居民，并保留到那一日，新郎对新娘说：\BibleRef{歌 4:7}\BibleQuote{我的爱卿！你全美丽，毫无瑕疵。}我们也应在这个意义上理解那话：\BibleRef{斐 2:15}\BibleQuote{好使你们成为无可指摘、纯洁无邪，天主无瑕的子女}；因为祂不是说\Emph{你们是}，而是说\Emph{你们应当是}。他是指着将来而言，并非陈述一个关于现在的实况；因此，这里有辛劳与努力，而在另一世界，则有劳苦与德行的赏报。最后，若望写道：\BibleRef{若一 3:2}\BibleQuote{亲爱的诸位，我们已是天主的子女，但我们将来如何，还没有显明；我们知道，当祂显现时，我们必要相似祂，因为我们要看见祂实在怎样。}那么，虽然我们是天主的子女，但相似天主以及真实地瞻仰天主，被应许给我们是在那时，当祂在威严中显现的时候。\par

14. 从这种膨胀的骄傲中，产生出祈祷中的鲁莽行为，就是你在致某位寡妇的信中关于圣人们应如何祈祷的指示所彰显的。你说：\Quote{祂理应向天主举起双手；他以良好的良心倾吐恳求，能说：}\InnerQuote{主，祢知道，我向祢伸开的双手是多么神圣、多么无辜、多么纯洁，远离一切欺骗、不义和抢夺；我向祢倾吐祈祷的嘴唇是多么正义、多么无瑕、远离一切谎言，求祢怜悯我。}这岂是基督徒的祈祷，还是像福音中那位骄傲的法利塞人（\BibleRef{路 18:11}）所说的：\BibleQuote{天主，我感谢祢，因为我不像其他人，勒索、不义、奸淫，也不像这个税吏：我每周两次禁食，凡我所得的，我都捐献十分之一。}他不过因天主的仁慈感谢天主，因为他不像其他人：他痛恨罪恶，并不将他的义德归为己有。但是你说：\Quote{如今祢知道，我在祢面前伸开的双手是多么神圣、多么无辜、多么纯洁，远离一切欺骗、不义和抢夺。}他说他每周两次禁食，为要折磨他邪恶放纵的肉身，并把一切所有的献上十分之一。因为\BibleRef{箴 13:8}：\BibleQuote{人的赎价是他的财富。}你与魔鬼一同夸口说：\Quote{我要升到星辰之上，我要将我的宝座安置在天上，我要与至高者相似。}达味说：\BibleQuote{我的腰充满幻象}；\BibleQuote{我的伤疮发臭流脓，由于我的愚昧}；\BibleQuote{求祢不要与祢的仆人进入审判}；\BibleQuote{在祢面前，凡活着的人没有一个是义的。}你夸耀自己是神圣、无辜、纯洁的，并向天主伸出洁净的双手。你不满足于夸耀你的一切行为，还要说你从一切言语的罪中洁净了；你还告诉我们，你的嘴唇是多么正义、无瑕、远离谎言。圣咏作者唱道：\BibleQuote{每一个人都是说谎者}；这有权威支持：圣保禄在\BibleRef{罗 3:4}中说：\BibleQuote{为使天主成为真实的，而每一个人都是说谎者}。而你却有正义、无瑕、远离谎言的嘴唇。依撒意亚哀叹说（\BibleRef{依 6:5}）：\BibleQuote{我有祸了！因为我完了！因为我是嘴唇不洁的人，住在嘴唇不洁的人民中间}；随后，一个色辣芬用钳子取来一块炽热的炭，洁净了先知的口，因为按你话语的指向，他并非骄傲，而是承认了自己的过犯。正如我们在圣咏中读到：\BibleQuote{诡诈的舌头，你当受什么报应？强者的利箭，加上毁灭性的炭火。}在这一切骄傲膨胀、祈祷中的夸耀和自以为圣洁之后，像一个愚妄的人试图说服另一个愚妄的人一样，你以这些话结束：\Quote{这些嘴唇，我用它们倾吐恳求，求祢怜悯我。}如果你圣洁，如果你无辜，如果你从一切污秽中得洁净，如果你在言语和行为上都没有犯罪——尽管雅各伯说（\BibleRef{雅 3:2}）：\BibleQuote{在言语上没有过失的，就是一个成全的人}，并说：\BibleQuote{没有人能制伏自己的舌头}——你怎么还求怜悯呢？因此，唉！你哀叹自己，并倾吐祈祷，因为你圣洁、纯洁、无辜，是一个嘴唇无瑕、远离谎言、赋有与天主相像的能力的人。基督在十字架上如此祈祷：\BibleQuote{我的天主，我的天主，祢为什么舍弃了我？为什么远离不救助我？}并且，又：如\BibleRef{路 23:46}所载：\BibleQuote{父啊！我把我的灵魂交托在祢手中}；又如\BibleRef{路 23:34}所载：\BibleQuote{父啊！宽赦他们吧！因为他们不知道他们做的是什么}。而这位就是祂——曾为我们谢恩，说过（\BibleRef{玛 11:25}）：\BibleQuote{父啊！天地的主宰！我称谢祢。}\par

15. 我们的主如此教导祂的宗徒，使得信徒每日在祭献祂身体的圣事中，大胆地说：\Quote{我们在天的父！愿你的名被尊为圣}；他们热切渴望天主的圣名——它本身已是圣的——在他们内被尊为圣；而你却说：\Quote{主，你知道，我的手是多么圣洁、多么无辜、多么纯净。} 然后他们说：\Quote{愿你的国来临}，盼望未来天国的希望，好使当基督为王时，罪绝不在他们必死的身体内作王；对此他们加上：\Quote{愿你的旨意奉行在人间，如同在天上}；如此，人的软弱得以效法天使，主的旨意得以在地上承行；而你却说：\Quote{人若愿意，就能完全无罪。} 宗徒祈求日用粮，或那超越一切食物的、将要来临的食粮，好使他们堪当领受基督的身体；而你却被你过度的圣德和稳固的义德所驱使，竟敢妄求天上的恩赐。接下来是：\Quote{宽免我们的罪债，如同我们也宽免得罪我们的人。} 他们刚从洗礼池中上来，藉着重生并合于我们的主救主，从而应验了关于他们所记载的：\Quote{那些过犯得蒙赦免、罪恶得蒙遮盖的人是有福的}，就在初领基督身体的圣事中，他们却说：\Quote{宽免我们的罪债}，尽管这些罪债在宣认基督时已蒙赦免；但你却傲慢自夸，夸耀你圣手的洁净和言语的纯正。无论一个人的悔改多么彻底，在犯罪和失败之后拥有德行多么完美，这样的人能像那些刚离开基督洗礼池的人一样毫无过犯吗？然而，这些后者奉命说：\Quote{宽免我们的罪债，如同我们也宽免得罪我们的人}；并非出于虚伪的谦卑，而是因为他们害怕人性的软弱，畏惧自己的良心。他们说：\Quote{不要让我们陷入诱惑}；而你和\Person{若维尼安}一致声称，那些怀着完整信心受洗的人不会再受诱惑或犯罪。最后，他们加上：\Quote{但救我们免于凶恶。} 为什么他们向主祈求他们已凭自由意志能力所拥有的事物？哦，人啊，你已在洗礼池中被洗净，论到你曾说过：\Quote{这位全身洁白、依靠她所爱者上升的是谁？} 因此，新娘被洗洁净，但她不能保持纯洁，除非蒙主扶持。你为何渴望天主的仁慈释放你——你刚刚才从罪恶中获得释放？唯一的解释是这一原则：我们坚持，当我们做完一切时，必须承认我们是无用的仆人。\par

16. 因此，你的祈祷胜过了法利塞人的骄傲，与税吏相比，你被定罪了。他远远站着，连举目望天都不敢，只是捶着胸膛说，\BibleRef{路 18:13}\Quote{天主，可怜我这个罪人吧！} 主的宣告正是基于此：\Quote{我告诉你们：这人下去，到他家里，成了正义的，而不是那个人。因为凡高举自己的，必被贬抑；凡贬抑自己的，必被高举。} 宗徒被贬抑，好使他们被高举。你的门徒被高举，好使他们跌倒。在你先前对那寡妇的奉承中，你竟不羞愧地说，如此在地上找到的虔诚和处处陌生的真理，在所有人中唯独与她同住。你不记得熟悉的经文：\BibleRef{依 3:12}\Quote{我的百姓啊！那引导你的使你走错路，并败坏你所行的路径。} 而你竟明言赞美她，说：\Quote{你是有福的，超乎一切想象！你是多么蒙福！如果义德——如今被认为只在天上——唯独在地上与你同在时才能找到。} 这是教导还是杀害？是将人从地上举起，还是从天上摔下，竟将天使不敢归于自身的事归于一个可怜的女子？如果虔诚、真理和义德在地上除了在一个女子身上别无寻处，那么我们在哪里能找到你那些自夸在地上无罪的义德追随者呢？这两章关于祈祷和赞美的话，你和你的门徒惯于发誓说不是你们的，然而在其中，你华丽的风格和西塞罗式措辞的优雅如此明显，尽管你像乌龟一样慢吞吞地昂首阔步，但你却不敢承认你在私下教导并公开贩卖的内容。幸福的人啊！你的著作除了你自家的门徒，无人抄写，这样，任何看来不可接受的，你都可以争辩说那不是你的作品，而是别人的。然而，谁有能力模仿你语言的魅力呢？\par

17. C. 我再也无法忍耐了；我的耐心完全被你的不义之言击败。请告诉我，恳求你，小婴儿们犯了什么罪？那些按\Person{约纳}先知所说不知左右手的人，既不能归咎于过犯的意识，也不能归咎于无知。他们不能犯罪，却可能丧亡；他们的膝盖太软弱，无法行走，他们发出含糊不清的哭声；我们嘲笑他们试图说话；而与此同时，可怜的不幸者！永恒的痛苦的折磨为他们准备好了。\par

A. 啊！现在你的门徒变成了老师，你开始变得流畅，甚至可以说是雄辩。\Person{安东尼}，一位卓越的演说家，\Person{西塞罗}大声赞扬他，说他见过许多口若悬河的人，但至今从未见过雄辩的演说家；所以，不要用那些并非来自你自家花园的修辞花朵来取悦我，那些花朵习惯于取悦缺乏经验和年少之人的耳朵，而是直截了当地告诉我你的想法。\par

C. 我要说的是——你至少必须承认，那些不能犯罪的人没有罪。\par

A. 我会承认这一点，如果他们已受洗归于基督；并且如果你不会立刻迫使我相信你的意见，即人若愿意，就能无罪；因为他们既无能力也无意愿；但他们藉着天主的恩宠，免于一切罪恶，这恩宠是他们在圣洗中所领受的。\par

C. 你逼我说出一句令人反感的话并问：为什麼，他们犯了什么罪？好叫你立刻在民众的骚动中用石头打死我。你没有能力杀我，但你确有这意愿。\par

A. 容许异端者成为异端者的人，就是杀害他。但我们责备他时，却赐给他生命；你可以死于你的异端，而活于公教信仰。\par

C. 如果你知道我们是异端者，为什麼不控告我们？\par

A. 因为《圣经》（\BibleRef{铎 3:10}）宗徒教导我，在第一次和第二次劝戒以后，要避开异端者，而非控告他。宗徒知道这样的人是乖僻的、自定己罪的。此外，让我的信仰取决于别人的判断，乃是极愚之举。因为假如有人称你为公教徒，我就立刻附和吗？谁为你辩护，说你的乖谬意见是正确的，他并不能救你脱离恶名，反而使自己背负不忠的罪名。你的众多支持者永远不能证明你是公教徒，只会显明你是异端者。然而，我愿意这类意见由教会的权威加以压制；否则，我们就像那些向啼哭的孩子展示可怕图画的人一样。愿对天主的敬畏赐予我们这件事——即轻视其他一切恐惧。因此，要么辩护你的意见，要么放弃你无法辩护的。无论谁被召来为你辩护，他必须被登记为党徒，而非保护人。\par

18. C. 请告诉我，并解除我的一切疑虑，为什麼要给孩子付洗？\par

A. 为使他们在圣洗中罪过得赦免。\par

C. 他们犯了什么罪？一个人若未被捆绑，怎能被释放？\par

A. 你问我！福音的号筒将回答，外邦人的导师、那照耀全世界的金器：\BibleRef{罗 5:14} \Quote{从亚当直到梅瑟，死亡作了王，连那些没有像亚当一样违背命令而犯罪的人，也属它权下；亚当原是那未来者的预像。}你若反驳说有些人被提及没有犯罪，你必须明白：他们不像亚当在乐园中违背天主的命令那样犯罪。但所有的人都要么因他们的祖先亚当，要么因他们自己而负责。婴儿在圣洗中从束缚其父亲的锁链中被释放。年龄大到有辨识能力的人，则因基督的血而从他自己的或别人的罪的锁链中被释放。你不可因我持此见解而视我为异端者，因为真福殉道者西彼廉（你自夸在《圣经》证明的分类上是他的对手）在写给费都斯主教关于婴儿圣洗的书信中这样说：\Quote{此外，如果连那些最恶劣的罪人，以及在受洗前大大冒犯天主的人，一旦相信，便获得罪过的赦免，没有人被阻止领受圣洗和恩宠，那麼，一个婴儿岂不更不该被阻止领受圣洗？因为他刚出生，没有犯过任何罪。他只是按肉身生在亚当的子孙中，由第一次出生而承担了古旧的死亡。而且，他更容易被允许获得罪赦，正是由于赦免给他的不是他自己的罪，而是别人的罪。因此，最亲爱的兄弟，我们在会议中决定：没有人该被我们阻止领受圣洗和天主的恩宠，祂对众人都是慈悲、温和、良善的。既然这条规则应对所有人遵守和持守，请记住，它更应对婴儿和刚出生的孩子本身遵守，因为他们更应得到我们的援助以获得天主的怜悯，因为他们的哭声和眼泪从出生起就是一项不断的祈祷。}\par

19. 那位圣洁的人、雄辩的主教奥古斯丁不久以前就婴儿圣洗写了二篇论文给马塞里努斯（此人后来虽无辜，却被异端者以参与赫拉克利安暴政的借口处死），以反对你们的异端，该异端主张：婴儿受洗不是为了罪过的赦免，而是为了得以进入天国，正如福音所载：\BibleRef{若 3:3} \Quote{人若不重生于水和圣神，不能进入天主的国。}他又写了第三篇给同一位马塞里努斯，反对那些像你们一样说人若愿意，没有天主的帮助也能无罪的人。而且，最近又写了第四篇给希拉里，反对你们这个充满乖谬的学说。据说他还有另外几篇特别针对你们的著作正在修订中，但尚未到我手中。因此，我想我必须放弃我的任务，免得荷拉斯的这句话落在我身上：\Quote{不要把木柴运进森林。}因为我们要么说与他同样的话，那便多余；要么若想说些新东西，便会发现最好的论点已被那卓越的天才抢先了。我有一件事要说，并以此结束我的论述：你们要么必须给我们一个新信经，好叫你们在把婴儿因父、及子、及圣神之名付洗之后，再为他们进入天国付洗；要么，如果对婴儿和对成年人是同一洗礼，那麼婴儿也当为罪过的赦免而受洗，正如亚当违命的相似。但你们若认为赦免别人的罪意味着不义，而那个不能犯罪的人不需要它，那么请转向你们特别喜爱的奥力振，他说在受洗时，很久以前在天上犯下的古老过犯得以解除。这样，你们不仅在其他事上受他的权威引导，在这件事上也将跟从他的错误。\par

\SourceRecord{Translated by W.H. Fremantle, G. Lewis and W.G. Martley.；Nicene and Post-Nicene Fathers, Second Series；Vol. 6.；Edited by Philip Schaff and Henry Wace.；Buffalo, NY: Christian Literature Publishing Co.,；1893.；<http://www.newadvent.org/fathers/30113.htm>.}
